\documentclass[11pt]{amsart}
\usepackage[lmargin=1in,rmargin=1in,tmargin=1in,bmargin=1in]{geometry}
\usepackage[ps,all,arc,rotate]{xy}
\usepackage{graphicx, float, epstopdf}
\usepackage{bbm}
\usepackage{color}
\usepackage[unicode,bookmarks=false]{hyperref}
\hypersetup{hidelinks}
\usepackage{centernot}
\usepackage{fancyhdr}
\usepackage{multirow}
\usepackage[utf8]{inputenc}
\usepackage{amsfonts,amssymb,amsmath,amsthm,mathrsfs}
\usepackage{graphics, setspace}
\usepackage{braket}
\usepackage{mathtools}
\usepackage{tikz}
\usepackage{upgreek}
\usepackage{xcolor}
\usepackage{array,esint}

\numberwithin{equation}{section}
\numberwithin{figure}{section}
\newcommand{\skipline}{\vspace{11pt}}  % might have to be preceeded by \\ in the previous line (clarity)
\allowdisplaybreaks[4]         %allows breaking multiline environments in amsmath commands
\newcommand{\blankline}{\quad\pagebreak[2]}
\newtheorem{lemma}{Lemma}[section]
\newtheorem{theorem}{Theorem}[section]
\newtheorem{proposition}{Proposition}[section]
\newtheorem{example}{Example}[section]
\newtheorem{claim}[lemma]{Claim}
\newtheorem{conjecture}[lemma]{Conjecture}
\newtheorem{corollary}[lemma]{Corollary}
\theoremstyle{definition}
\newtheorem{definition}{Definition}[section]
\newtheorem{remark}{Remark}[section]
\usepackage{booktabs}
\usepackage{array}
\usepackage{vcell}
\newcommand{\Z}{\mathbb{Z}}
\newcommand{\Q}{\mathbb{Q}}
\newcommand{\R}{\mathbb{R}}
\newcommand{\C}{\mathbb{C}}
\newcommand{\N}{\mathbb{N}}
\newcommand{\Nset}{\mathcal{N}}
\newcommand{\Mset}{\mathcal{M}}
\newcommand{\Oset}{\mathcal{O}}
\newcommand{\M}{\mathcal{M}}
\newcommand{\B}{\mathcal{B}}
\newcommand{\Hset}{\mathcal{H}}
\newcommand{\Pri}{\mathbb{P}}
\newcommand{\partition}{p}
\newcommand{\F}{\mathcal{F}}
\newcommand{\ID}{\mathbf{1}}
\newcommand{\id}{\textrm{id}}
\newcommand{\Ha}{{\hat H}}
\newcommand{\card}{\operatorname{card}}
\newcommand{\vt}{{\vec \theta }}
\newcommand{\I}{\textnormal{i}}
\newcommand{\e}{\operatorname{e}}
\newcommand{\Li}{\operatorname{Li}}
\newcommand{\real}{\operatorname{Re}}
\newcommand{\imag}{\operatorname{Im}}
\newcommand{\BS}{\operatorname{BS}}
\newcommand{\SV}{\operatorname{SV}}
\newcommand{\SVH}{\operatorname{SVH}}
\newcommand{\SVA}{\operatorname{SVA}}
\newcommand{\sgn}{\operatorname{sgn}}
\newcommand{\sumtwo}{\operatorname*{\sum\sum}}
\newcommand{\sumthree}{\operatorname*{\sum\sum\sum}}
\newcommand{\sumdots}{\operatorname*{\sum\cdots\sum}}
\newcommand{\sumdotss}{\operatorname*{\sum\sum\cdots\sum}}
\newcommand{\ointdots}{\operatorname*{\oint\cdots\oint}}
\newcommand{\intdots}{\operatorname*{\int_2^\infty \cdots \int_2^\infty}}
\newcommand{\modu}{\operatorname{mod}}
\def\sumstar{\operatornamewithlimits{\sum\nolimits^*}}
\newcommand{\quadand}{\quad \textnormal{and} \quad}
\newcommand{\distinct}{\operatorname{distinct}}
\newcommand{\barg}{{\bar g}}
\newcommand{\barr}{{\bar r}}
\newcommand{\barK}{{\bar K}}
\newcommand{\barq}{{\bar q}}
\newcommand{\bark}{{\bar k}}
\newcommand{\barh}{{\bar h}}
\newcommand{\barrr}{{\bar r}}
\newcommand{\barL}{{\bar L}}
\newcommand{\barl}{{\bar l}}
\newcommand{\tr}{\operatorname{tr}}
\newcommand{\norm}[1]{\left\lVert#1\right\rVert}
\newcommand{\threeprimes}{p_1^{k_1}p_2^{k_2} \cdots p_r^{k_r}}
\newcommand{\threeells}{\ell_1^{k_1}\ell_2^{k_2} \cdots \ell_r^{k_r}}
\newcommand{\imod}[1]{\ (\mathrm{mod}\ #1)}

\renewcommand{\S}{\mathfrak{S}}
\renewcommand{\M}{\mathfrak{M}}
\newcommand{\m}{\mathfrak{m}}


\newcommand{\AGK}{A_{\mathrm{GK}}}
\newcommand{\Aarc}{A_{\mathrm{arc}}}
\newcommand{\AIng}{A_{\mathrm{Ing}}}
\newcommand{\Res}{\operatorname*{Res}}

\newcommand{\ttriple}[3]{\textcolor{#3}{(#1,#2)}}

% Manuscript-specific macros.
\newcommand{\eps}{\varepsilon}
\newcommand{\abs}[1]{\lvert #1\rvert}
\newcommand{\Fp}{\mathbb F_p}
\newcommand{\PP}{\mathbb P}
\newcommand{\Qc}{\mathcal Q}
\newcommand{\rad}{\operatorname{rad}}
\newcommand{\adj}{\operatorname{adj}}
\emergencystretch=2em


\title[Character sums over zeros of quadratic forms]{Character sums over zeros of quaternary quadratic forms}
\author[H.~Nath]{Hemjyoti Nath}
\address{Department of Mathematics, University of Florida, Gainesville, FL 32611, USA}
\email{h.nath@ufl.edu}
\author[N.~Robles]{Nicolas Robles}
\address{Quantum Sky, Project 137, 11950 Democracy Drive, Suite 650, Reston, VA 20190, USA}
\email{robles.nicolas.m@gmail.com}
\author[A.~Zaharescu]{Alexandru Zaharescu}
\address{Department of Mathematics, University of Illinois Urbana-Champaign, Altgeld Hall, 1409 W. Green Street, Urbana, IL 61801, USA}
\email{zaharesc@illinois.edu}
\subjclass[2020]{Primary 11D09; Secondary 11L40, 11N36, 11P55, 14G05}
\keywords{Quaternary quadratic forms, character sums, circle method, delta symbol, square sieve, power sieve}

\begin{document}

\begin{abstract}
Let $F(x)=x^{\mathsf T}Ax$ be a quadratic form in four variables,
where $A$ is a symmetric integral matrix whose determinant is neither
zero nor a square, and let $G$ be a form of even degree with integer
coefficients such that the curve $F=G=0$ is smooth. We
prove an asymptotic formula for the sum of the Jacobi symbols
$(G(x)/q)$ over the integral zeros $x$ of $F$, with a smooth weight
$W(x/B)$, for squarefree $q$ coprime to $2\det A$ and to the primes of
bad reduction of the curve. The error term is
$(Bq)^{1+\varepsilon}$, uniformly in $q$ and in the coefficients of
$G$, and the local factors of the main term are character sums over
the quadric modulo the primes dividing $q$. The same is proved for
products of characters of order $r$ modulo the primes dividing $q$,
when $r$ divides the degree of $G$. As an
application, the number of zeros $x$ of $F$ with $\lvert x_i\rvert\le B$
at which $G$ is a nonzero square or $r$th power, for such $r$, is $O(B^{5/3+\varepsilon}+B^{1+\varepsilon}\log^2H)$, where
$H\ge2$ bounds the coefficients of $G$.
\end{abstract}

\maketitle

\section{Introduction}\label{sec:intro}

\subsection{Previous results in the literature}\label{sec:literature}

Let $A$ be a nonsingular symmetric $4\times4$ matrix with integer
entries, and set
\begin{equation}\label{eq:forms}
 F(x)=x^{\mathsf T}Ax,\qquad \alpha=\det A,\qquad N_F=2\abs{\alpha}.
\end{equation}
We assume throughout that $\alpha$ is not a square. Every quadratic
form in four variables with integer coefficients is of this shape
after multiplication by $2$, which does not change its zeros. Let
$\chi$ be the primitive quadratic character attached to the field
$\Q(\sqrt\alpha)$, so that $\chi(p)=(\alpha/p)$ for every prime
$p\nmid N_F$, where $(\cdot/p)$ is the Legendre symbol.

The zeros of $F$ in a region of size $B$ are counted, with a weight
$W\in C_c^\infty(\R^4)$, by the sum of $W(x/B)$ over the $x\in\Z^4$
with $F(x)=0$. Heath-Brown \cite[Theorem~6]{HB} proved, for a class
of weights and by means of a form of the delta symbol of Duke,
Friedlander and Iwaniec \cite{DFI}, that
\begin{equation}\label{eq:HB}
 \sum_{\substack{x\in\Z^4\\F(x)=0}}W(x/B)
 =B^2\sigma_\infty(W,F)\,\S_F+O_{F,W,\eps}(B^{3/2+\eps}).
\end{equation}
Here the singular integral and the local densities of $F$ are
\begin{equation}\label{eq:singint}
 \sigma_\infty(W,F)=\int_{F(y)=0}\frac{W(y)}{\abs{\nabla F(y)}}\,dS(y),
 \qquad
 \sigma_p=\lim_{k\to\infty}p^{-3k}\#\{x\bmod p^k:F(x)\equiv0\imod{p^k}\},
\end{equation}
where $dS$ is the surface measure on the cone $F(y)=0$ in $\R^4$, and
the singular series is
\begin{equation}\label{eq:SF}
 \S_F=L(1,\chi)\prod_p\Big(1-\frac{\chi(p)}p\Big)\sigma_p.
\end{equation}
The integral and the product converge absolutely, and the limits
exist; see Sections~\ref{sec:delta} and~\ref{sec:moduli}. Getz
\cite[Theorem~1.1]{G} improved the error term in \eqref{eq:HB} to
$O(B^{1+\eps})$.

The present paper concerns the sums over the zeros of $F$ in which
the summand carries a character. For an odd squarefree integer
$q\ge1$ we write $\chi_q(n)=(n/q)$ for the Jacobi symbol. Thus
$\chi_q=\prod_{p\mid q}\chi_p$, $\chi_q(n)=0$ when $(n,q)>1$, and
$\chi_1(n)=1$ for every $n$. For a form $G$ of even degree with
integer coefficients, and $B\ge2$, set
\begin{equation}\label{eq:Tdef}
 T_{q,G}(B)=\sum_{\substack{x\in\Z^4\\F(x)=0}}W(x/B)\chi_q(G(x)).
\end{equation}
These sums were studied by Iwaniec and Munshi \cite{IM} for the
diagonal quadratic forms
\begin{equation}\label{eq:diagonal}
 F(x)=\sum_{i=1}^4a_ix_i^2,\qquad G(x)=\sum_{i=1}^4b_ix_i^2,
 \qquad a_ib_j-a_jb_i\ne0\quad(1\le i<j\le4),
\end{equation}
for which $\alpha=a_1a_2a_3a_4$. Their aim was a bound for
\begin{equation}\label{eq:MG}
 M_G(B)=\#\{x\in\Z^4:\lVert x\rVert_\infty\le B,\ F(x)=0,\
 G(x)\text{ is a nonzero square}\},
\end{equation}
where $\lVert x\rVert_\infty=\max_i\abs{x_i}$; as explained in
\cite{IM}, it controls the number of rational points of height at
most $B$, outside the lines, on the quartic del Pezzo surface
$F(x)=0$, $y^2=G(x)$. By the square
sieve of Heath-Brown \cite{HB2}, $M_G(B)$ is bounded in terms of the
sums $T_{pp',G}(B)$ for pairs of primes $p\ne p'$, and
\cite[Proposition~2.1]{IM} is the bound
\begin{equation}\label{eq:IM}
 T_{q,G}(B)\ll\Big(\frac{B^2}q+B^{3/2}+Bq^2\Big)(Bq)^\eps
\end{equation}
for squarefree $q$, from which $M_G(B)\ll B^{9/5+\eps}$ follows
\cite[Theorem~1.1]{IM}. Mallesham \cite{M} applied the delta symbol
to the equation $F(x)/q=0$ on the vectors $x$ with
$F(x)\equiv0\imod q$, as in the work of Browning and Munshi
\cite{BM}, and obtained $M_G(B)\ll B^{5/3+\eps}$
\cite[Theorem~1.3]{M}.

Both proofs use bounds for sums over the zeros of $F$ modulo a
prime, namely \cite[Lemma~4.6]{IM} and \cite[Lemmas~5.3 and~5.4]{M},
which do not hold at all frequencies; see
Remark~\ref{rem:literature}. The bound \eqref{eq:IM} is nevertheless
correct when $q$ is prime to $2\alpha$ and to the numbers
$a_ib_j-a_jb_i$: it follows from Theorem~\ref{thm:polynomial} below,
since $Bq\le Bq^2$. The corresponding bound
\cite[Proposition~2.1]{M} has the term $B^2q^{-3/2}$ in place of
$B^2/q$, whereas $T_{q,G}(B)$ has exact order $B^2/q$ for the primes
$q$ satisfying the same condition, with $\chi(q)=-1$ and
$B^\eta\le q\le B^{1/2-\eta}$, when
$\sigma_\infty(W,F)\S_F\ne0$ and $\eta>0$ is fixed; this follows from
Theorem~\ref{thm:diagonal} and Corollary~\ref{cor:relative}. The two
bounds for $M_G(B)$ are contained in
Theorem~\ref{thm:uniformcount}.

For a fixed modulus $q$ the sums \eqref{eq:Tdef} are covered by a
theorem of Getz \cite[Theorem~3.3]{G}, on the zeros of quadratic forms
weighted by an arbitrary function at the finite places, which gives
an asymptotic formula with error term $O(B^{1+\eps})$ and an implied
constant that depends on $q$ in an unspecified way. Lindqvist
\cite{L} studied the zeros of $F$ in a residue class to a fixed
modulus, where a further term of order $B^2$ can occur, and Huang,
Schindler and Shute \cite{HSS} obtained asymptotic formulas for a
residue class to a modulus that grows with $B$. The error term of
\cite[Theorem~4.1]{HSS} is for one residue class; when it is added
over the residue classes modulo $q$ on which $F$ vanishes, the result
is larger than the error term of Theorem~\ref{thm:polynomial} below
by a power of $q$. Sums over the zeros
of $F$ with arbitrary coefficients attached to one of the
coordinates are considered in \cite{GSS}.

The number $M_G(B)$ can also be approached by the determinant
method: for a quadratic form $G$, a theorem of Salberger
\cite[Theorem~0.6]{S} on intersections of two quadrics gives the
exponent $3/2$, with a power of the size of the coefficients of $G$;
see Section~\ref{sec:fixed}. Heath-Brown and Pierce \cite{HBP} used
the sieve with characters of order $r$ to count rational points on
cyclic covers of projective space, and Bucur, Cojocaru, Lal\'in and
Pierce \cite{BCLP} gave geometric generalizations of the square
sieve. Polynomial sieves were applied by Bonolis and Pierce \cite{BP}
and by Bonolis, Pierce and Woo \cite{BPW} to thin sets of type II in
affine space, in \cite{BPW} with explicit dependence on the
coefficients, through the uniform stratification of Bonolis,
Kowalski and Woo \cite{BKW}. Browning and Loughran \cite{BL}
developed sieves for rational points on quadrics.

\subsection{New results}\label{sec:results}

Let $G\in\Z[x_1,x_2,x_3,x_4]$ be a form of degree $2d$, where $d\ge1$,
and let $H_G$ be the maximum of $2$ and of the absolute values of the
coefficients of $G$. We assume that
\begin{equation}\label{eq:smooth}
 \operatorname{rank}\big(\nabla F(x),\nabla G(x)\big)=2
 \qquad\text{for every }x\in\overline{\Q}^4\setminus\{0\}
 \text{ with }F(x)=G(x)=0.
\end{equation}
Equivalently, the equations $F=G=0$ define a smooth complete
intersection curve in $\PP^3$, the intersection of the quadric
surface $F=0$ with the hypersurface $G=0$. For a prime $p\nmid N_F$,
we say that $G$ has good reduction at $p$ if \eqref{eq:smooth} holds
for the reductions of $F$ and $G$ modulo $p$, with $\overline{\Fp}$ in
place of $\overline{\Q}$, and that it has bad reduction at $p$
otherwise. Set
\begin{equation}\label{eq:NG}
 N_G=\rad(N_F)\prod_{\mathclap{\substack{p\nmid N_F\\G\text{ has bad reduction at }p}}}\;p,
\end{equation}
where $\rad(n)$ is the product of the primes dividing $n$. It is part
of Theorem~\ref{thm:polynomial} that this product is finite. The
primes dividing $N_G$ are those excluded from the modulus $q$.

We study the sums \eqref{eq:Tdef} for such forms $G$ and for
squarefree $q$ prime to $N_G$. Their main term involves the
quantities \eqref{eq:singint} and \eqref{eq:SF} and the following
local factors.

For a prime $p\nmid N_F$, let $\Qc(\Fp)$ be the set of points of
$\PP^3(\Fp)$ at which $F$ vanishes. It has $p^2+(1+\chi(p))p+1$
elements, by Lemma~\ref{lem:gauss}. For $p\nmid N_G$ set
\begin{equation}\label{eq:apdef}
 a_p(G)=\sum_{P\in\Qc(\Fp)}\chi_p(G(P)).
\end{equation}
The summand does not depend on the coordinates chosen for $P$, because
$G$ has even degree. A point $P$ lies below $1+\chi_p(G(P))$ points of
the double cover $y^2=G(x)$ of the quadric, in which $y$ has weight
$d$, so that $a_p(G)$ is the
difference between the numbers of points over $\Fp$ of the cover and
of the quadric.

\begin{theorem}\label{thm:polynomial}
Let $d\ge1$. There is a constant $C_{F,d}>0$ such that
\begin{equation}\label{eq:badheight}
 N_G\ll_{F,d}H_G^{C_{F,d}}
\end{equation}
for every form $G$ of degree $2d$ satisfying \eqref{eq:smooth}. For
every such $G$, every squarefree $q\ge1$ with $(q,N_G)=1$, and
$B\ge2$,
\begin{equation}\label{eq:mainthm}
 T_{q,G}(B)=B^2\sigma_\infty(W,F)\S_G(q)
 +O_{F,d,W,\eps}\big((Bq)^{1+\eps}\big),
\end{equation}
where
\begin{equation}\label{eq:Sgeneral}
 \S_G(q)=\S_F\prod_{p\mid q}\frac{a_p(G)(p^2-1)}{p^2\,\#\Qc(\Fp)}
 =\S_F\prod_{p\mid q}\frac{a_p(G)(p-1)(p-\chi(p))}{p^2(p^2-\chi(p))}.
\end{equation}
Moreover $\abs{a_p(G)}\le(8d^2-8d+4)p$ for every prime $p\nmid N_G$,
so that
\begin{equation}\label{eq:mainbound}
 \abs{\S_G(q)}\ll_{F,d,\eps}q^{-1+\eps},
 \qquad
 T_{q,G}(B)\ll_{F,d,W,\eps}\Big(\frac{B^2}q+Bq\Big)(Bq)^\eps.
\end{equation}
The implied constants do not depend on the coefficients of $G$.
\end{theorem}

The factor at $p$ in \eqref{eq:Sgeneral} is $1-p^{-2}$ times the mean
value of $\chi_p(G(P))$ over $\Qc(\Fp)$; the factor $1-p^{-2}$ comes
from the vectors $x$ divisible by $p$, at which the character
vanishes. The bound for $T_{q,G}(B)$ in \eqref{eq:mainbound} is
smaller than the trivial bound $O(B^2)$ by a power of $B$ when
$B^\eta\le q\le B^{1-\eta}$ for a fixed $\eta>0$. For $q=1$ the
formula \eqref{eq:mainthm} is \eqref{eq:HB} with the error term of
\cite[Theorem~1.1]{G}. For a fixed $q$, an asymptotic formula with
error term $O_q(B^{1+\eps})$ follows from \cite[Theorem~3.3]{G}; what
\eqref{eq:mainthm} adds is the error term uniform in $q$ and in the
coefficients of $G$ together, and the evaluation of the local
factors.

As an application we bound the number $M_G(B)$ of \eqref{eq:MG}, of
zeros of $F$ at which $G$ is a nonzero square.

\begin{theorem}\label{thm:uniformcount}
Let $d\ge1$, and let $G$ be a form of degree $2d$ satisfying
\eqref{eq:smooth}. For $B\ge2$,
\begin{equation}\label{eq:uniformcount}
 M_G(B)\ll_{F,d,\eps}B^{5/3+\eps}+B^{1+\eps}\log^2(2H_GB).
\end{equation}
The implied constant does not depend on the coefficients of $G$.
\end{theorem}

In particular, for every fixed $K_0>0$, one has
$M_G(B)\ll_{F,d,K_0,\eps}B^{5/3+\eps}$ uniformly for $H_G\le B^{K_0}$. The
proof applies the square sieve of Heath-Brown \cite{HB2}, with
Theorem~\ref{thm:polynomial} at the moduli $q=pp'$. For a quadratic
form $G$, a theorem of Salberger \cite{S} gives the bound
$M_G(B)\ll_\eps(H_G^{1/2}B)^{3/2+\eps}$, which is at least as strong
as \eqref{eq:uniformcount} when $H_G\le B^{2/9}$; see
Section~\ref{sec:fixed}.

When $F$ and $G$ are the diagonal quadratic forms of
\eqref{eq:diagonal}, the local factors can be evaluated. Set
\begin{equation}\label{eq:kappadef}
 \kappa_i=-a_i\prod_{j\ne i}(a_ib_j-a_jb_i)\quad(1\le i\le4),
 \qquad M=\prod_{1\le i<j\le4}(a_ib_j-a_jb_i).
\end{equation}

\begin{theorem}\label{thm:diagonal}
Let $F$ and $G$ be as in \eqref{eq:diagonal}. Then
$N_G=\rad(2\alpha M)$, and for every prime $p\nmid2\alpha M$,
\begin{equation}\label{eq:diagtrace}
 a_p(G)=p\sum_{i=1}^4\chi_p(\kappa_i).
\end{equation}
Consequently, for every squarefree $q\ge1$ with $(q,2\alpha M)=1$,
the formula \eqref{eq:mainthm} holds with
\begin{equation}\label{eq:diagmain}
 \S_G(q)=\S_F\prod_{p\mid q}\bigg(\frac{(p-1)(p-\chi(p))}{p(p^2-\chi(p))}
 \sum_{i=1}^4\chi_p(\kappa_i)\bigg).
\end{equation}
\end{theorem}

Since $\kappa_1\kappa_2\kappa_3\kappa_4=\alpha M^2$, the sum in
\eqref{eq:diagtrace} is $\pm2$ when $\chi(p)=-1$, and is $0$ or
$\pm4$ when $\chi(p)=1$.

The method applies to characters of higher order. Let $r\ge3$ and
$d\ge1$, and let $G$ now be a form of degree $rd$ satisfying
\eqref{eq:smooth}; the quantities $H_G$ and $N_G$ are defined as
before. For each prime $p\equiv1\imod r$ with $p\nmid N_G$, let
$\psi_p$ be a character modulo $p$ of exact order $r$. For a
squarefree $q\ge1$ all of whose prime divisors are of this kind, set
$\psi_q=\prod_{p\mid q}\psi_p$, with $\psi_1(n)=1$ for every $n$.
Since $r$ divides the degree of $G$, the value $\psi_p(G(P))$ does not
depend on the coordinates chosen for a point $P\in\Qc(\Fp)$, and we
set
\begin{equation}\label{eq:apsidef}
 a_{p,\psi}(G)=\sum_{P\in\Qc(\Fp)}\psi_p(G(P)).
\end{equation}

\begin{theorem}\label{thm:cyclic}
Let $r\ge3$ and $d\ge1$. There is a constant $C_{F,r,d}>0$ such that
$N_G\ll_{F,r,d}H_G^{C_{F,r,d}}$ for every form $G$ of degree $rd$
satisfying \eqref{eq:smooth}. For every such $G$, every $q$ and
$\psi_q$ as above, and $B\ge2$,
\begin{equation}\label{eq:cyclicmain}
 \sum_{\substack{x\in\Z^4\\F(x)=0}}W(x/B)\psi_q(G(x))
 =B^2\sigma_\infty(W,F)\S_{G,\psi}(q)
 +O_{F,r,d,W,\eps}\big((Bq)^{1+\eps}\big),
\end{equation}
where
\begin{equation}\label{eq:cyclicfactor}
 \S_{G,\psi}(q)=\S_F\prod_{p\mid q}
 \frac{a_{p,\psi}(G)(p^2-1)}{p^2\,\#\Qc(\Fp)}.
\end{equation}
Moreover
$\abs{a_{p,\psi}(G)}\le2\big(1+(rd-1)^2\big)p$. The implied
constants do not depend on the coefficients of $G$ or on the choice
of the characters $\psi_p$.
\end{theorem}

\begin{corollary}\label{cor:powercount}
Let $r\ge3$ and $d\ge1$, let $G$ be a form of degree $rd$ satisfying
\eqref{eq:smooth}, and let
\[
 M_{G,r}(B)=\#\{x\in\Z^4:\lVert x\rVert_\infty\le B,\ F(x)=0,\
 G(x)=k^r\text{ for some integer }k\ne0\}.
\]
For $B\ge2$,
\begin{equation}\label{eq:powercount}
 M_{G,r}(B)\ll_{F,r,d,\eps}B^{5/3+\eps}+B^{1+\eps}\log^2(2H_GB).
\end{equation}
The implied constant does not depend on the coefficients of $G$.
\end{corollary}

For even $r$ this follows from Theorem~\ref{thm:uniformcount}, since
an $r$th power is a square. The new cases are those of odd $r$, for
instance the cube values of a cubic form on the zeros of $F$. When
$G$ has degree $r$, the determinant method gives
$M_{G,r}(B)\ll_{r,\eps}(H_G^{1/r}B)^{m_r+\eps}$ with
$m_r=\max\{1,3/\sqrt{2r}\}$, which has a smaller power of $B$ than
\eqref{eq:powercount} unless $H_G$ exceeds a fixed power of $B$; see
Section~\ref{sec:fixed}.

Finally, suppose that $\sigma_\infty(W,F)\S_F\ne0$ and that
$a_p(G)\ne0$ for every $p\mid q$. Then \eqref{eq:mainthm} is an
asymptotic formula uniformly for $q\le B^{1/2-\eta}$ when $G$ is
quadratic, and for $q\le B^{1/3-\eta}$ in general, for every fixed
$\eta>0$. For the characters of Theorem~\ref{thm:cyclic}, if
$\sigma_\infty(W,F)\S_F\ne0$ and $a_{p,\psi}(G)\ne0$ for every
$p\mid q$, the corresponding range is
$q\le B^{2/(\varphi(r)+4)-\eta}$, where $\varphi$ is Euler's
function. These statements are proved in
Section~\ref{sec:relative}. Section~\ref{sec:character-sums} also
contains an estimate for the sum \eqref{eq:Tdef} with $\chi_q(G(x))$
replaced by $\prod_{p\mid q}(\chi_p(G(x))-\beta_p)$, where $\beta_p$
is the factor at $p$ in \eqref{eq:Sgeneral}.

\subsection{Outline of the proofs}\label{sec:outline}

The theorems of Section~\ref{sec:results} are deduced from an
asymptotic formula for the sums of $W(x/B)\xi_q(x)$ over the zeros
$x$ of $F$, where $q$ is squarefree and $\xi_q$ is the product over
the primes $p\mid q$ of
functions $\xi_p$ on $\Fp^4$ that vanish at the origin, are bounded
by $1$ and are invariant under multiplication of $x$ by nonzero
scalars (Proposition~\ref{prop:analytic}). The formula is proved
under one hypothesis on the functions $\xi_p$, which concerns the
sums $C_p(u)$ of \eqref{eq:Cdef}, the Fourier transform of $\xi_p$
on the zeros of $F$ modulo $p$: they are $O(p^{3/2})$, except at the
zeros $u$ of a
nonzero polynomial $D_p$ of bounded degree, where they are $O(p^2)$.
The polynomial $D_p$ may depend on $p$.

As in \cite{M}, the condition $F(x)\equiv0\imod q$ is imposed first,
and the delta symbol is applied to the integer $F(x)/q$, with the
parameter $Q=B/\sqrt q$. The Poisson summation formula modulo $qc$
then expresses the sum as a sum over the moduli $c\ll Q$ and the
frequencies $u\in\Z^4$ of the products of the complete sums
$\mathcal S_{q,c}(u)$ of \eqref{eq:Snewdef} and of the integrals of
\cite{HB}; the
frequencies that matter have length about $\sqrt q$
(Lemma~\ref{lem:poisson}). The sums $\mathcal S_{q,c}(u)$ factor
over the primes dividing $qc$. At a prime dividing both $q$ and $c$
they are
evaluated in closed form, in terms of the value of $\xi_p$ at
$\adj(A)u$ and of a Ramanujan sum of $F^*(u)$; here $\adj(A)$ is the
adjugate of $A$ and $F^*(u)=u^{\mathsf T}\adj(A)u$
(Lemmas~\ref{lem:ramified} and~\ref{lem:goodglobal}).

The frequencies are of three kinds. The frequency $u=0$ gives the
main term; the singular series converges by means of the character
$\chi$, and the factor at $p\mid q$ is
$C_p(0)p^{-3}$ divided by the density $\sigma_p$
(Lemmas~\ref{lem:singseries} and~\ref{lem:zerofreq}). For the
frequencies with $F^*(u)\ne0$ absolute values are taken. The sum over
$c$ then costs a factor $(Bq)^\eps$, and there remains an average of
$\abs{C_{q_2}(u)}$ over the $u$ with $q_1\mid F^*(u)$, for the
factorizations $q=q_1q_2$. In this average the larger bound for
$C_p(u)$ on the zeros of $D_p$ is harmless, because these $u$ are
counted by a lemma on polynomial congruences in a box that is
uniform in the polynomials (Lemmas~\ref{lem:polynomialbox},
\ref{lem:genericmod} and~\ref{lem:average}). For the frequencies
$u\ne0$ with $F^*(u)=0$ the sum over $c$ is carried out first. It
has cancellation because $\chi$ is not principal, that is, because
$\alpha$ is not a square, and the number of these frequencies is
bounded by Lemma~\ref{lem:quadriccount}
(Lemmas~\ref{lem:conemod} and~\ref{lem:oncone}). The three
contributions give the error term $(Bq)^{1+\eps}$.

In Section~\ref{sec:fourier} the hypothesis is verified for
$\xi_p(x)=\psi(G(x))$, where $\psi$ is a nonprincipal character
modulo $p$ whose order divides the degree of $G$. The sum $C_p(0)$ is $p-1$ times the
sum of $\psi(G(P))$ over the points $P$ of the quadric $F=0$ over
$\Fp$, which is bounded by a theorem of Katz \cite{K}
(Lemma~\ref{lem:katz}). For $p\nmid u$, the sum $C_p(u)$ is
expressed by this sum and by the sum over the plane section
$u\cdot x=0$ of the quadric. When $F^*(u)\ne0$ this section is a conic, which is
parametrized by the projective line, and the bound of Weil applies
if the plane is transversal to the curve $F=G=0$
(Lemmas~\ref{lem:conic} and~\ref{lem:section}). The planes that are
not transversal are among the zeros of a polynomial obtained by
elimination, whose degree is bounded in terms of $F$ and of the
degree of $G$ only; the same elimination bounds $N_G$ by a power of
$H_G$ (Lemma~\ref{lem:exceptional}). For
diagonal quadratic forms the sums $C_p(u)$ are computed for every
$u$ (Lemma~\ref{lem:pencil}).

Theorem~\ref{thm:uniformcount} and Corollary~\ref{cor:powercount}
follow from the square sieve and its analogue for $r$th powers, with
the bound for the character sums at the moduli $q=pp'$
(Lemma~\ref{lem:sieve}). The ranges of $q$ at the end of
Section~\ref{sec:results} come from lower bounds for the local
factors that do not vanish: a nonzero $a_p(G)$ is an integer, which
is divisible by $p$ when $G$ is a quadratic form, and for characters
of order $r\ge3$ the norm of a nonzero $a_{p,\psi}(G)$ from the field
of $r$th roots of unity is a nonzero integer
(Lemmas~\ref{lem:divisible} and~\ref{lem:lowerbeta}).

\subsection{Notation}\label{sec:notation}

Vectors $x$, $y$ and $u$ have four components; $x\cdot u$ is the
standard inner product, $\lVert x\rVert$ the Euclidean norm, and
$\lVert x\rVert_\infty=\max_i\abs{x_i}$. We write
$\e(t)=\exp(2\pi it)$ and $\e_c(t)=\e(t/c)$ for an integer $c\ge1$.
The letter $p$ denotes a prime, $v_p(n)$ is the exponent of $p$ in
$n$, and $(m,n)$ is the greatest common divisor. As usual, $\varphi$,
$\mu$, $\tau$ and $\omega$ denote Euler's function, the M\"obius
function, the number of divisors and the number of distinct prime
divisors, and $\rad(n)$ is the product of the primes dividing $n$. A
star on a sum over $a\bmod c$ restricts it to $(a,c)=1$, and
$\overline a$ denotes an inverse of $a$ to a modulus that is clear
from the context. For a prime $p$, $\Fp$ is the field with $p$
elements and $\overline{\Fp}$ an algebraic closure of it; a character
modulo $p$ is extended to $\Fp$ by the value $0$ at $0$. The symbol
$\mathbf 1_{\mathcal E}$ is $1$ when the condition $\mathcal E$ holds
and $0$ otherwise. The letter $\eps$ denotes an arbitrarily small
positive number, not necessarily the same at each occurrence. Implied
constants depend only on the parameters shown as subscripts; in
particular they do not depend on the coefficients of $G$ unless $G$
is a subscript.

\subsection{Organization of the paper}\label{sec:organization}

Section~\ref{sec:lemmas} contains the lemmas that are used
throughout: the delta symbol and the estimates of \cite{HB} for its
integrals, a bound for the number of zeros of a quaternary quadratic
form and a lemma on polynomial congruences, and the bounds of Weil and Katz
for character sums over finite fields. Section~\ref{sec:local}
studies the complete sums $\mathcal S_{q,c}(u)$, their sums over the
modulus $c$, the singular series, and the averages over the
exceptional frequencies. In Section~\ref{sec:fourier} the
hypotheses of Section~\ref{sec:local} are verified for the
functions $\psi(G(x))$, and the sums $C_p(u)$ are evaluated for
diagonal quadratic forms.
Section~\ref{sec:proofs} contains the proofs: the asymptotic formula
for general weights in Section~\ref{sec:analytic},
Theorems~\ref{thm:polynomial}, \ref{thm:diagonal}
and~\ref{thm:cyclic} in Section~\ref{sec:character-sums},
Theorem~\ref{thm:uniformcount} and Corollary~\ref{cor:powercount} in
Section~\ref{sec:sieve}, and the ranges of $q$ in which
\eqref{eq:mainthm} is an asymptotic formula in
Section~\ref{sec:relative}. Section~\ref{sec:fixed} compares
Theorem~\ref{thm:uniformcount} and Corollary~\ref{cor:powercount}
with the determinant method, and Section~\ref{sec:conclusion}
contains some remarks and questions for future work.

\section{Preliminary lemmas}\label{sec:lemmas}

\subsection{The delta method and its integrals}\label{sec:delta}

We detect the equation $F(x)=0$ by the delta symbol of Duke,
Friedlander and Iwaniec \cite{DFI}, in the form given by Heath-Brown
\cite{HB}.

\begin{lemma}\label{lem:delta}
There are an infinitely differentiable function $h$ on
$(0,\infty)\times\R$ and, for every $Q\ge3$, a constant $c_Q>0$, such
that for every integer $n$
\begin{equation}\label{eq:delta}
 \mathbf 1_{n=0}=\frac{c_Q}{Q^2}\sum_{c\ge1}\ \sum_{a\bmod c}^{*}
 \e_c(an)\,h\Big(\frac cQ,\frac n{Q^2}\Big).
\end{equation}
Moreover $c_Q=1+O_j(Q^{-j})$ for every $j\ge1$, one has
$h(x,y)\ll x^{-1}$ for all $x$ and $y$, and $h(x,y)=0$ unless
$x\le\max\{1,2\abs y\}$.
\end{lemma}

\begin{proof}
This is \cite[Theorem~1]{HB}, where it is stated for $Q>1$. The
function $h$ is defined in \cite[Section~3]{HB} and does not depend
on $Q$. The constant $c_Q$ is defined by \cite[(3.3)]{HB}, which
requires an integer strictly between $Q/2$ and $Q$; there is one for
$Q>2$, but not for $Q=2$.
\end{proof}

By the last assertion of the lemma, the sum over $c$ in
\eqref{eq:delta} is finite.

We turn to the integrals that arise when \eqref{eq:delta} is combined
with the Poisson summation formula, and to the estimates of \cite{HB}
for them. By a weight we mean an infinitely differentiable function
on $\R^4$ with compact support. For a weight $w$, a real number
$r>0$ and $v\in\R^4$ set
\begin{equation}\label{eq:Istar}
 I_r^*(v;w)=\int_{\R^4}w(x)\,h(r,F(x))\,
 \e\Big(-\frac{v\cdot x}r\Big)\,dx.
\end{equation}
This is the integral $I_r^*(v)$ of \cite[Section~7]{HB}. The
estimates of \cite{HB} are proved for weights of small support, and
a general weight is decomposed as follows. Let
\[
 \theta(x)=c_0^{-4}\prod_{i=1}^4w_0(x_i),\qquad
 w_0(t)=\begin{cases}\exp\big(-(1-t^2)^{-1}\big),&\abs t<1,\\
 0,&\abs t\ge1,\end{cases}
\]
where $c_0$ is the integral of $w_0$ over $\R$. Thus $\theta$ is a
nonnegative weight supported in $[-1,1]^4$, with integral $1$. For a
weight $W$, a number $0<\delta\le1$ and $y\in\R^4$ set
\begin{equation}\label{eq:pieces}
 W_y(x)=W(x)\,\theta\Big(\frac{x-y}\delta\Big),
 \qquad\text{so that}\qquad
 W(x)=\delta^{-4}\int_{\R^4}W_y(x)\,dy.
\end{equation}
The weight $W_y$ vanishes identically unless $y$ lies within distance
$2\delta$ of the support of $W$.

\begin{lemma}\label{lem:integrals}
Let $0<\rho_1<\rho_2$, and let $\mathcal W$ be a set of weights
supported in $\rho_1\le\lVert x\rVert\le\rho_2$, whose partial
derivatives of each order are bounded on $\R^4$ uniformly in
$\mathcal W$. There are numbers $\delta\in(0,1]$ and $r_0\ge1$,
depending only on $F$ and $\mathcal W$, with the following
properties. Let $W\in\mathcal W$ and $y\in\R^4$, and let $w=W_y$ be
as in \eqref{eq:pieces}. Then $I_r^*(v;w)=0$ for $r\ge r_0$ and every
$v\in\R^4$. Moreover, either $F(x)\ne0$ for every $x$ in the support
of $w$, or, for $r>0$,
\begin{align}
 I_r^*(0;w)&=\sigma_\infty(w,F)+O_j(r^j),
 \qquad\frac{\partial I_r^*(0;w)}{\partial r}\ll r^{-1},\label{eq:Izero}\\
 I_r^*(v;w)&\ll_jr^{-1}\lVert v\rVert^{-j},\label{eq:Irapid}\\
 \abs{I_r^*(v;w)}+r\,\bigg|\frac{\partial I_r^*(v;w)}{\partial r}\bigg|
 &\ll_\eps r^{1-2\eps}\lVert v\rVert^{-1+\eps},\label{eq:Idecay}
\end{align}
for every integer $j\ge1$, every $0<\eps<1/2$, and, in
\eqref{eq:Irapid} and \eqref{eq:Idecay}, every $v\ne0$. The
derivatives are taken with $v$ fixed, and the implied constants
depend only on $F$, $\mathcal W$ and the parameters shown.
\end{lemma}

\begin{proof}
Let $r_0=1+2\max\abs{F(x)}$, the maximum being over
$\lVert x\rVert\le\rho_2$. For $r\ge r_0$ the function $h(r,F(x))$
vanishes on the support of $w$, by Lemma~\ref{lem:delta}, and
$I_r^*(v;w)=0$.

The other assertions restate results of \cite{HB}, which are proved
there for the classes of weights $\mathcal C(S)$ and $\mathcal C_0(S)$
of \cite[Sections~2 and~6]{HB}; here $S$ is a set of parameters, on
which the implied constants may depend. We apply them with four
variables and with the form $F$, and let $S$ consist of $F$,
$\mathcal W$ and the number
\[
 \lambda=\min\{\abs{\nabla F(x)}:\rho_1\le\lVert x\rVert\le\rho_2\},
\]
which is positive because $\nabla F(x)=2Ax$ and $A$ is nonsingular.
Then $\mathcal W\subseteq\mathcal C(S)$, and $\abs{\nabla F}\ge\lambda$
on the support of every $W\in\mathcal W$, as is required in
\cite[Section~1]{HB}. The functions $W_y$ are those attached to $W$
in \cite[Lemma~2]{HB}. By \cite[Lemma~12]{HB} there is a $\delta$,
bounded below in terms of $S$, such that for each $y$ either
$\abs{F(x)}\gg_S1$ on the support of $W_y$, or $W_y$ belongs to
$\mathcal C_0(S)$ after a suitable permutation of the variables. As
the proof of that lemma shows, one of the variables may also have to
be replaced by its negative. In the second case there is therefore a
linear map $\iota$ of $\R^4$, which permutes the coordinates and
changes the sign of at most one of them, such that $W_y\circ\iota$
belongs to the class $\mathcal C_0(S)$ formed with the form
$F\circ\iota$. Since $\iota$ is orthogonal, the substitution
$x\mapsto\iota x$ in \eqref{eq:Istar} and in \eqref{eq:singint} shows
that $I_r^*(v;W_y)$ and $\sigma_\infty(W_y,F)$ are equal to the
corresponding quantities for the weight $W_y\circ\iota$, the form
$F\circ\iota$ and the vector $\iota^{-1}v$, which has the same norm
as $v$. There are finitely many maps $\iota$, and each $F\circ\iota$
is a nonsingular quadratic form. Hence it suffices to prove
\eqref{eq:Izero}, \eqref{eq:Irapid} and \eqref{eq:Idecay} for a
weight $w\in\mathcal C_0(S)$.

For such a weight the surface $F=0$ is, on the support of $w$, the
graph of a function $x_1=x_1(x_2,x_3,x_4)$, on which
$dS=\abs{\nabla F}(\partial F/\partial x_1)^{-1}dx_2\,dx_3\,dx_4$.
Thus $\sigma_\infty(w,F)$ is the singular integral of
\cite[Theorem~3]{HB}. The first part of \eqref{eq:Izero} is
\cite[Lemma~13]{HB} for $r\ll_S1$; for the other values of $r$ it
holds because both terms are $O_S(1)$: the integral $I_r^*(0;w)$ by
the bound $h(x,y)\ll x^{-1}$, and $\sigma_\infty(w,F)$ by the formula
just given. By \cite[Lemma~14]{HB},
\[
 I_r^*(v;w)\ll r^{-1}\abs{I(r;r^{-1}v)},\qquad
 \frac{\partial I_r^*(v;w)}{\partial r}\ll r^{-2}\abs{I(r;r^{-1}v)},
\]
where $I(r;u)$ stands for integrals of the kind defined there, not
necessarily the same in the two bounds; the second bound uses that
$F$ is homogeneous. For these integrals,
\cite[Lemmas~15, 18 and~22]{HB} give
\[
 I(r;u)\ll_Sr,\qquad
 I(r;u)\ll_{S,j}r^{-j}\lVert u\rVert^{-j},\qquad
 I(r;u)\ll_{S,\eps}(r^{-1}\lVert u\rVert)^\eps r\lVert u\rVert^{-1},
\]
the last because $F$ is a nonsingular quadratic form in four
variables. With $u=r^{-1}v$ these are the second part of
\eqref{eq:Izero}, then \eqref{eq:Irapid}, and \eqref{eq:Idecay}.
\end{proof}

A weight $W$ whose support does not contain the origin is covered by
Lemma~\ref{lem:integrals} with $\mathcal W=\{W\}$. We also record
that the integral defining $\sigma_\infty(W,F)$ in
\eqref{eq:singint} converges absolutely for every weight $W$: the
cone $F=0$ is smooth away from the origin, on it
$\abs{\nabla F(y)}=2\abs{Ay}\gg_F\lVert y\rVert$, and by homogeneity
its part with $\rho<\lVert y\rVert\le2\rho$ has surface measure
$O_F(\rho^3)$.

\subsection{Zeros of quaternary forms and polynomial congruences}\label{sec:congruences}

The first lemma bounds the number of zeros of a quaternary quadratic
form in a box.

\begin{lemma}\label{lem:quadriccount}
Let $J(z)=z^{\mathsf T}Sz$, where $S$ is a nonsingular symmetric
$4\times4$ matrix with rational entries whose determinant is not the
square of a rational number. Then, for $U\ge1$,
\begin{equation}\label{eq:quadriccount}
 \#\{z\in\Z^4:\lVert z\rVert\le U,\ J(z)=0\}\ll_{J,\eps}U^{2+\eps}.
\end{equation}
\end{lemma}

\begin{proof}
If $J$ has no zero in $\Q^4$ other than the origin, there is nothing
to prove. Otherwise let $v\in\Q^4\setminus\{0\}$ with $J(v)=0$, and
write $b(x,y)=x^{\mathsf T}Sy$, so that
$J(x+y)=J(x)+2b(x,y)+J(y)$. Since $S$ is nonsingular, there is
$w\in\Q^4$ with $b(v,w)\ne0$. Replacing $w$ by $w+tv$ with
$t=-J(w)/(2b(v,w))$, and then by a multiple of itself, we may assume
that $J(w)=0$ and $b(v,w)=1/2$. The vectors $y\in\Q^4$ with
$b(v,y)=b(w,y)=0$ form a subspace of dimension $2$, which meets the
span of $v$ and $w$ only in the origin. Let $v_3,v_4$ be a basis of
it, and let $T$ be the matrix with columns $v,w,v_3,v_4$. Then $T$ is
invertible and
\[
 J(Ty)=y_1y_2+H_0(y_3,y_4)\qquad(y\in\Q^4),
\]
where $H_0$ is a binary quadratic form with rational coefficients.
Comparing determinants gives $\det(S)\det(T)^2=D/16$, where $D$ is the
discriminant of $H_0$. If $H_0$ had a zero in $\Q^2$ other than the
origin, then $D$ would be the square of a rational number, and so
would $\det S$. Hence $H_0(y_3,y_4)\ne0$ unless $y_3=y_4=0$.

Choose an integer $k\ge1$ such that $kT^{-1}$ has integer entries and
$H=kH_0$ has integer coefficients. For $z\in\Z^4$ with
$\lVert z\rVert\le U$ and $J(z)=0$, the vector $y=kT^{-1}z$ has
integer components, determines $z$, and satisfies
$\lVert y\rVert\ll_JU$ and
\[
 ky_1y_2+H(y_3,y_4)=0.
\]
If $(y_3,y_4)\ne(0,0)$, then $n=H(y_3,y_4)$ is a nonzero integer with
$\abs n\ll_JU^2$, and $y_1$ divides $n$. By the divisor bound there
are $O_{J,\eps}(U^\eps)$ possibilities for $y_1$, each of which
determines $y_2$. These $y$ number $O_{J,\eps}(U^{2+\eps})$. If
$y_3=y_4=0$, then $y_1y_2=0$, and there are $O_J(U)$ such $y$.
\end{proof}

Let $\adj(A)$ be the adjugate of $A$, so that $A\adj(A)=\alpha I$,
and set
\begin{equation}\label{eq:dualform}
 F^*(u)=u^{\mathsf T}\adj(A)u.
\end{equation}
The determinant of $\adj(A)$ is $\alpha^3$, which is not a square, so
that Lemma~\ref{lem:quadriccount} applies to $F$ and to $F^*$. By a
dyadic decomposition it gives, for $U\ge1$,
\begin{equation}\label{eq:conecount}
 \sum_{\substack{0<\lVert u\rVert\le U\\F^*(u)=0}}\frac1{\lVert u\rVert}
 \ll_{F,\eps}U^{1+\eps}.
\end{equation}

The second lemma counts the solutions of polynomial congruences to a
squarefree modulus in a box. The bound does not depend on the
coefficients of the polynomials, which may be different at each
prime.

\begin{lemma}\label{lem:polynomialbox}
Let $n,E\ge1$ be integers and let $m\ge1$ be squarefree. For every
prime $p\mid m$ let $P_p\in\Fp[z_1,\ldots,z_n]$ be a nonzero
polynomial of total degree at most $E$. Let $I_1,\ldots,I_n$ be
intervals, each containing at most $Y$ integers, where $Y\ge1$. Then
\begin{equation}\label{eq:polynomialbox}
 \#\{z\in\Z^n:z_i\in I_i\ (1\le i\le n),\ P_p(z)=0\text{ in }\Fp
 \text{ for every }p\mid m\}
 \le C_{n,E}^{\omega(m)}\Big(\frac{Y^n}m+Y^{n-1}\Big),
\end{equation}
where $C_{n,E}=E+E^2+\cdots+E^n$.
\end{lemma}

\begin{proof}
For $m=1$ the count is at most $Y^n$. Let $m>1$ and $k=\omega(m)$. We
argue by induction on $n$.

Let $n=1$. A nonzero polynomial in $\Fp[z_1]$ of degree at most $E$
has at most $E$ zeros in $\Fp$. By the Chinese remainder theorem, the
integers counted lie in at most $E^k$ residue classes modulo $m$, and
each class contains at most $Y/m+1$ integers of $I_1$. This gives the
bound $E^k(Y/m+1)$.

Let $n\ge2$, and write $z=(t,y)$ with $t\in\Z$ and $y\in\Z^{n-1}$.
For each $p\mid m$ write
\[
 P_p(t,y)=\sum_{j=0}^EP_{p,j}(y)t^j,\qquad P_{p,j}\in\Fp[z_2,\ldots,z_n].
\]
For $y\in\Z^{n-1}$, let $g(y)$ be the product of the primes $p\mid m$
for which $P_{p,j}(y)=0$ in $\Fp$ for every $j$. At a prime dividing
$m/g(y)$ the polynomial $P_p(t,y)$ in $t$ is nonzero. By the case
$n=1$ with the modulus $m/g(y)$, the number of $t\in I_1$ for which
$(t,y)$ is counted in \eqref{eq:polynomialbox} is at most
\[
 E^k\Big(\frac{Yg(y)}m+1\Big).
\]
We sum this over the $y$ with $y_i\in I_{i+1}$ for $1\le i\le n-1$.
Since $g(y)$ is the sum of $\varphi(h)$ over the divisors $h$ of
$g(y)$,
\[
 \sum_yg(y)=\sum_{h\mid m}\varphi(h)\#\{y:h\mid g(y)\}.
\]
Let $h\mid m$. For each $p\mid h$ choose an index $j(p)$ such that
$P_{p,j(p)}$ is not the zero polynomial; there is one, because
$P_p\ne0$. If $h\mid g(y)$, then $P_{p,j(p)}(y)=0$ in $\Fp$ for
every $p\mid h$. With $C=C_{n-1,E}$, the induction hypothesis for the
modulus $h$ and the polynomials $P_{p,j(p)}$ gives
\[
 \#\{y:h\mid g(y)\}\le C^{\omega(h)}\Big(\frac{Y^{n-1}}h+Y^{n-2}\Big).
\]
The number of $y$ is at most $Y^{n-1}$. Hence the count in
\eqref{eq:polynomialbox} is at most
\[
 E^k\bigg(\frac{Y^n}m\sum_{h\mid m}\frac{\varphi(h)}hC^{\omega(h)}
 +\frac{Y^{n-1}}m\sum_{h\mid m}\varphi(h)C^{\omega(h)}+Y^{n-1}\bigg).
\]
Here
\[
 \sum_{h\mid m}\frac{\varphi(h)}hC^{\omega(h)}
 =\prod_{p\mid m}\Big(1+C\,\frac{p-1}p\Big)\le(1+C)^k,
 \qquad
 \sum_{h\mid m}\varphi(h)C^{\omega(h)}
 =\prod_{p\mid m}\big(1+C(p-1)\big)\le C^km,
\]
and $C^k+1\le(1+C)^k$ because $k\ge1$. The count is therefore at
most $(E(1+C))^k(Y^n/m+Y^{n-1})$, and $E(1+C_{n-1,E})=C_{n,E}$.
\end{proof}

Since $C_{n,E}^{\omega(m)}\ll_{n,E,\eps}m^\eps$, the right side of
\eqref{eq:polynomialbox} is $\ll_{n,E,\eps}m^\eps(Y^n/m+Y^{n-1})$.

\subsection{Bounds over finite fields}\label{sec:finite-fields}

We use two bounds for sums of a character modulo $p$. The first is
Weil's bound for polynomials in one variable.

\begin{lemma}\label{lem:weil}
Let $\psi$ be a nonprincipal character modulo $p$, and let
$g\in\Fp[z]$ be a polynomial of degree $m\ge1$ with $m$ distinct
zeros in $\overline{\Fp}$. Then
\begin{equation}\label{eq:weil}
 \bigg|\sum_{z\in\Fp}\psi(g(z))\bigg|\le(m-1)\sqrt p.
\end{equation}
\end{lemma}

\begin{proof}
This is \cite[(1.2)]{K}. It is also a case of
\cite[Theorem~11.23]{IK}, since a polynomial of positive degree whose
zeros are simple is not a constant multiple of a perfect power.
\end{proof}

We shall apply it to sums over the projective line, in the following
form.

\begin{corollary}\label{cor:weilline}
Let $n\ge2$, let $\psi$ be a nonprincipal character modulo $p$ such
that $\psi^n$ is principal, and let $g\in\Fp[s,t]$ be a form of
degree $n$ with $n$ distinct zeros in $\PP^1(\overline{\Fp})$. Then
\begin{equation}\label{eq:weilline}
 \bigg|\sum_{(s:t)\in\PP^1(\Fp)}\psi(g(s,t))\bigg|\le(n-1)\sqrt p+1.
\end{equation}
\end{corollary}

\begin{proof}
The summand does not depend on the coordinates chosen for the point,
because $\psi(\lambda^n)=1$ for $\lambda\in\Fp\setminus\{0\}$. The
points of $\PP^1(\Fp)$ are $(1:0)$ and the points $(z:1)$ with
$z\in\Fp$. The zeros of the polynomial $g(z,1)$ in $\overline{\Fp}$
are simple, since they correspond to zeros of $g$. Its degree is $n$
if $g(1,0)\ne0$. Otherwise $t$ divides $g$ and $t^2$ does not, so
that the degree is $n-1\ge1$. By Lemma~\ref{lem:weil} the sum over
the points $(z:1)$ is at most $(n-1)\sqrt p$ in absolute value, and
the point $(1:0)$ contributes at most $1$.
\end{proof}

The second bound is for sums over the quadric $F=0$. For a prime
$p\nmid N_F$ let $\Qc(\Fp)$ be the set of points of $\PP^3(\Fp)$ at
which $F$ vanishes, as in Section~\ref{sec:results}.

\begin{lemma}\label{lem:katz}
Let $p\nmid N_F$ be a prime and $n\ge2$. Let $\psi$ be a nonprincipal
character modulo $p$ such that $\psi^n$ is principal, and let
$G\in\Fp[x_1,x_2,x_3,x_4]$ be a form of degree $n$ such that
\begin{equation}\label{eq:smoothp}
 \operatorname{rank}\big(\nabla F(x),\nabla G(x)\big)=2
 \qquad\text{for every }x\in\overline{\Fp}^4\setminus\{0\}
 \text{ with }F(x)=G(x)=0.
\end{equation}
Then
\begin{equation}\label{eq:katz}
 \bigg|\sum_{P\in\Qc(\Fp)}\psi(G(P))\bigg|\le(2n^2-4n+4)p.
\end{equation}
\end{lemma}

\begin{proof}
Let $\Qc$ be the surface $F=0$ in $\PP^3$ over $\Fp$. Since
$p\nmid N_F$, it is smooth. It is therefore geometrically
irreducible, since two components of a surface in $\PP^3$ meet, and
the surface is singular where they meet. By \eqref{eq:smoothp} the
hypersurface $G=0$ is transverse to $\Qc$. By
\cite[Theorem~8.1(1)]{K}, the cohomology groups with compact supports of
the complement of $G=0$ in $\Qc$ over $\overline{\Fp}$, with
coefficients in the sheaf attached to $\psi(G)$, vanish in every
degree other than $2$. In degree $2$ the group has dimension
\[
 C_0=\int_{\Qc}\frac{c(\Qc)}{1+nL},
\]
where $c(\Qc)$ is the total Chern class of $\Qc$ and $L$ is the class
of a plane section, and the eigenvalues of the Frobenius of $\Fp$ on
it have absolute value $p$. By the Lefschetz trace formula, as in
\cite[(5.3)]{K}, the sum in \eqref{eq:katz} is the trace of this
Frobenius. Hence it is at
most $C_0p$ in absolute value. For a quadric surface,
$c(\Qc)=(1+L)^4(1+2L)^{-1}=1+2L+2L^2$ and $\int_{\Qc}L^2=2$, so that
\[
 C_0=\int_{\Qc}(1+2L+2L^2)(1-nL+n^2L^2)=2(n^2-2n+2).\qedhere
\]
\end{proof}

With $n=2d$ and $\psi=\chi_p$ the right side of \eqref{eq:katz} is
$(8d^2-8d+4)p$, as in Theorem~\ref{thm:polynomial}, and with $n=rd$
it is the bound for $a_{p,\psi}(G)$ in Theorem~\ref{thm:cyclic}.

\section{Local sums}\label{sec:local}

Throughout this section, $N\ge1$ is an integer divisible by
$\rad(N_F)$, and $E,K\ge1$ are fixed. For every prime $p\nmid N$ we
are given a function $\xi_p:\Fp^4\to\C$ and a nonzero polynomial
$D_p\in\Fp[u_1,u_2,u_3,u_4]$ of total degree at most $E$. We assume
that
\begin{equation}\label{eq:weightproperties}
 \xi_p(0)=0,\qquad\abs{\xi_p(x)}\le1,\qquad\xi_p(tx)=\xi_p(x)
 \qquad(x\in\Fp^4,\ t\in\Fp\setminus\{0\}).
\end{equation}
For $x\in\Z^4$ we write $\xi_p(x)$ for the value of $\xi_p$ at the
reduction of $x$ modulo $p$. For a squarefree integer $q\ge1$ with
$(q,N)=1$ we set $\xi_q(x)=\prod_{p\mid q}\xi_p(x)$, with
$\xi_1(x)=1$, and
\begin{equation}\label{eq:Cdef}
 C_q(u)=\sum_{\substack{x\bmod q\\F(x)\equiv0\imod q}}\xi_q(x)\e_q(u\cdot x)
 \qquad(u\in\Z^4).
\end{equation}
We assume further that, for every prime $p\nmid N$ and every
$u\in\Z^4$,
\begin{equation}\label{eq:Cstrata}
 \abs{C_p(u)}\le K(p^{3/2}+p^2\,\mathbf 1_{D_p(u)=0}),
\end{equation}
where $D_p(u)=0$ is understood in $\Fp$. The results of this section
are applied in Section~\ref{sec:proofs} with $\xi_p(x)=\chi_p(G(x))$
and with $\xi_p(x)=\psi_p(G(x))$, and with $N=N_G$; the hypothesis
\eqref{eq:Cstrata} is verified for these functions in
Section~\ref{sec:fourier}. It is used only for the bound
\eqref{eq:Sxibound}, in Section~\ref{sec:average} and in
Section~\ref{sec:analytic}.

\subsection{Factorization and Gauss sums}\label{sec:gauss}

For integers $c\ge1$, squarefree $q\ge1$ with $(q,N)=1$, and
$u\in\Z^4$, set
\begin{align}
 S_c(u)&=\sum_{a\bmod c}^{*}\ \sum_{x\bmod c}\e_c(aF(x)+u\cdot x),\label{eq:Sdef}\\
 \mathcal S_{q,c}(u)&=\sum_{a\bmod c}^{*}\
 \sum_{\substack{x\bmod qc\\F(x)\equiv0\imod q}}
 \xi_q(x)\e_{qc}(aF(x)+u\cdot x).\label{eq:Snewdef}
\end{align}
The condition $F(x)\equiv0\imod q$ makes the summand in
\eqref{eq:Snewdef} a function of $a$ modulo $c$. Thus
$\mathcal S_{1,c}=S_c$ and $\mathcal S_{q,1}=C_q$. The sums
$\mathcal S_{q,c}(u)$ arise in Section~\ref{sec:analytic}, where the
delta method is applied to the equation $F(x)/q=0$ on the vectors
$x$ with $F(x)\equiv0\imod q$. We also write
\begin{equation}\label{eq:ramanujan}
 R_n(t)=\sum_{a\bmod n}^{*}\e_n(at)=\sum_{h\mid(n,t)}h\mu(n/h)
\end{equation}
for the Ramanujan sum; see \cite[Theorem~4.1]{MV}.

\begin{lemma}\label{lem:crt}
Let $q=q'q''$ be squarefree with $(q,N)=1$, and let $c=c'c''$ with
$(q'c',q''c'')=1$. Then
\begin{equation}\label{eq:crt}
 \mathcal S_{q,c}(u)=\mathcal S_{q',c'}(u)\mathcal S_{q'',c''}(u).
\end{equation}
Moreover $\mathcal S_{q,c}(tu)=\mathcal S_{q,c}(u)$ for every integer
$t$ with $(t,qc)=1$.
\end{lemma}

\begin{proof}
Set $n'=q'c'$ and $n''=q''c''$. As $x'$ runs over $(\Z/n'\Z)^4$ and
$x''$ over $(\Z/n''\Z)^4$, the vector $x=n''x'+n'x''$ runs over
$(\Z/qc\Z)^4$. Since $n''$ is invertible modulo $q'$ and $n'$ modulo
$q''$, the invariance in \eqref{eq:weightproperties} gives
$\xi_q(x)=\xi_{q'}(x')\xi_{q''}(x'')$, and the condition
$F(x)\equiv0\imod q$ is equivalent to $F(x')\equiv0\imod{q'}$ and
$F(x'')\equiv0\imod{q''}$. Moreover
\[
 F(x)\equiv n''^2F(x')+n'^2F(x'')\imod{n'n''},
 \qquad
 \frac{u\cdot x}{n'n''}=\frac{u\cdot x'}{n'}+\frac{u\cdot x''}{n''}.
\]
Similarly $a=c''a'+c'a''$ runs over the reduced residues modulo $c$
as $a'$ and $a''$ run over those modulo $c'$ and $c''$. Since
$F(x')/q'$ and $F(x'')/q''$ are integers, it follows that
\[
 \e_{qc}(aF(x)+u\cdot x)
 =\e_{n'}(c''n''a'F(x')+u\cdot x')\,\e_{n''}(c'n'a''F(x'')+u\cdot x'').
\]
Since $c''n''$ is invertible modulo $c'$ and $c'n'$ modulo $c''$, we
may replace $c''n''a'$ by $a'$ and $c'n'a''$ by $a''$. This proves
\eqref{eq:crt}.

For the last assertion, let $\overline t$ be an inverse of $t$ modulo
$qc$. In $\mathcal S_{q,c}(tu)$ replace $x$ by $\overline tx$ and
then $a$ by $t^2a$. By \eqref{eq:weightproperties} the weight and the
condition $F(x)\equiv0\imod q$ are unchanged, and the phase becomes
$aF(x)+u\cdot x$.
\end{proof}

In particular $C_q(u)=\prod_{p\mid q}C_p(u)$, and $S_c(u)$ is
multiplicative in $c$.

\begin{lemma}\label{lem:gauss}
Let $p\nmid N_F$ be a prime and $k\ge1$.
\begin{enumerate}
\item[\textup{(i)}] For every integer $a$ with $p\nmid a$,
\begin{equation}\label{eq:gauss}
 \sum_{x\bmod p^k}\e_{p^k}(aF(x)+u\cdot x)
 =\chi(p)^kp^{2k}\e_{p^k}\big(-\overline{4a\alpha}\,F^*(u)\big),
\end{equation}
where the inverse is taken modulo $p^k$.
\item[\textup{(ii)}] One has
\begin{equation}\label{eq:unram}
 S_{p^k}(u)=\chi(p)^kp^{2k}R_{p^k}(F^*(u)).
\end{equation}
\item[\textup{(iii)}] The number of points of $\PP^3(\Fp)$ at which
$F$ vanishes is $\#\Qc(\Fp)=p^2+(1+\chi(p))p+1$.
\end{enumerate}
\end{lemma}

\begin{proof}
(i) Since $p\nmid2a\alpha$, the matrix $A$ is invertible modulo $p^k$,
with inverse $\overline\alpha\adj(A)$. Replacing $x$ by
$x-\overline{2a\alpha}\adj(A)u$ changes $aF(x)+u\cdot x$ into
$aF(x)-\overline{4a\alpha}F^*(u)$ modulo $p^k$. It remains to show
that
\begin{equation}\label{eq:gauss0}
 V_k(a)=\sum_{x\bmod p^k}\e_{p^k}(aF(x))=\chi(p)^kp^{2k}.
\end{equation}
Let $k\ge2$, and write $x=y+p^{k-1}z$ with $y$ modulo $p^{k-1}$ and
$z$ modulo $p$. Then
$F(x)\equiv F(y)+2p^{k-1}y^{\mathsf T}Az\imod{p^k}$, and the sum over
$z$ of $\e_p(2ay^{\mathsf T}Az)$ is $p^4$ if $p\mid y$ and is $0$
otherwise. Writing $y=pw$ with $w$ modulo $p^{k-2}$, we obtain
$V_k(a)=p^4V_{k-2}(a)$, where $V_0(a)=1$. For $k=1$, a linear change
of variables over $\Fp$ brings $F$ to a diagonal form
$\alpha_1y_1^2+\cdots+\alpha_4y_4^2$, where
$\alpha_1\alpha_2\alpha_3\alpha_4$ is the product of $\alpha$ and a
nonzero square. For $b\in\Fp\setminus\{0\}$,
\[
 \sum_{y\in\Fp}\e_p(by^2)=\sum_{n\in\Fp}\big(1+\chi_p(n)\big)\e_p(bn)
 =\chi_p(b)\tau_p,\qquad\tau_p=\sum_{n\in\Fp}\chi_p(n)\e_p(n).
\]
Here $\overline{\tau_p}=\chi_p(-1)\tau_p$ and $\abs{\tau_p}^2=p$ by
\cite[Theorems~9.5 and~9.7]{MV}, so that $\tau_p^4=p^2$. Hence
\[
 V_1(a)=\chi_p(a^4\alpha_1\alpha_2\alpha_3\alpha_4)\tau_p^4=\chi(p)p^2.
\]
Together with the recursion this proves \eqref{eq:gauss0} for every
$k\ge1$.

(ii) Sum \eqref{eq:gauss} over the reduced residues $a$ modulo $p^k$,
and note that $-\overline{4a\alpha}$ runs over the reduced residues
with $a$.

(iii) By \eqref{eq:gauss0} with $k=1$, the number of $x\in\Fp^4$ with
$F(x)=0$ is
\[
 \frac1p\sum_{a\bmod p}V_1(a)=\frac1p\big(p^4+(p-1)\chi(p)p^2\big)
 =p^3+\chi(p)(p^2-p).
\]
Removing the origin and dividing by $p-1$ gives the assertion.
\end{proof}

At the primes dividing $N_F$ we use the following bound.

\begin{lemma}\label{lem:badprime}
For every integer $c\ge1$ and every $u\in\Z^4$,
\begin{equation}\label{eq:badbound}
 \abs{S_c(u)}\le4\abs\alpha^{1/2}c^3.
\end{equation}
\end{lemma}

\begin{proof}
Let $(a,c)=1$ and
$V=\sum_{x\bmod c}\e_c(aF(x)+u\cdot x)$. Writing $\abs V^2$ as a sum
over $x$ and $x+z$, we get
\[
 \abs V^2=\sum_{z\bmod c}\e_c(aF(z)+u\cdot z)
 \sum_{x\bmod c}\e_c(2ax^{\mathsf T}Az)
 \le c^4\#\{z\bmod c:2Az\equiv0\imod c\}.
\]
Let $d_1,\ldots,d_4$ be the invariant factors of the integer matrix
$2A$, so that $d_1d_2d_3d_4=\abs{\det(2A)}=16\abs\alpha$. The number
of $z$ modulo $c$ with $2Az\equiv0\imod c$ is
$\prod_i(d_i,c)\le16\abs\alpha$. Hence $\abs V\le4\abs\alpha^{1/2}c^2$,
and \eqref{eq:badbound} follows on summing over $a$.
\end{proof}

\subsection{Moduli not coprime to \texorpdfstring{$q$}{q}}\label{sec:ramified}

We now evaluate $\mathcal S_{q,c}(u)$ when $c$ is divisible by primes
dividing $q$.

\begin{lemma}\label{lem:ramified}
Let $p\nmid N$ be a prime and $k\ge1$. Then
\begin{equation}\label{eq:ramified}
 \mathcal S_{p,p^k}(u)
 =\chi(p)^{k+1}p^{2k+1}\,\xi_p(\adj(A)u)\,R_{p^{k+1}}(F^*(u)).
\end{equation}
\end{lemma}

\begin{proof}
For $j\ge2$ let
\[
 \mathcal U_j(u)=\sum_{a\bmod p^j}^{*}\ \sum_{x\bmod p^j}
 \xi_p(x)\e_{p^j}(aF(x)+u\cdot x).
\]
We first show that
\begin{equation}\label{eq:Uj}
 \mathcal U_j(u)=\chi(p)^jp^{2j}\,\xi_p(\adj(A)u)\,R_{p^j}(F^*(u)).
\end{equation}
Fix $a$ with $p\nmid a$, and write $x=y+p^{j-1}z$ with $y$ modulo
$p^{j-1}$ and $z$ modulo $p$. Then $\xi_p(x)=\xi_p(y)$, and
\[
 aF(x)+u\cdot x\equiv aF(y)+u\cdot y+p^{j-1}(2aAy+u)\cdot z\imod{p^j},
\]
because $2j-2\ge j$. The sum over $z$ is $p^4$ if
$2aAy+u\equiv0\imod p$, and is $0$ otherwise. The congruence means
that $y\equiv-\overline{2a\alpha}\adj(A)u\imod p$, and for these $y$
we have $\xi_p(y)=\xi_p(\adj(A)u)$ by \eqref{eq:weightproperties}.
The inner sum in $\mathcal U_j(u)$ is therefore
\[
 \xi_p(\adj(A)u)\,p^4
 \sum_{\substack{y\bmod p^{j-1}\\2aAy+u\equiv0\imod p}}
 \e_{p^j}(aF(y)+u\cdot y).
\]
The same computation without the weight shows that $p^4$ times the
last sum is the left side of \eqref{eq:gauss} with $k=j$. Summing
\eqref{eq:gauss} over $a$ as in Lemma~\ref{lem:gauss}(ii) gives
\eqref{eq:Uj}.

Now let $k\ge1$. The reduced residues modulo $p^{k+1}$ are the
numbers $a+tp^k$, where $a$ runs over the reduced residues modulo
$p^k$ and $t$ over the residues modulo $p$. The sum over $t$ of
$\e_p(tF(x))$ is $p$ if $p\mid F(x)$, and is $0$ otherwise. Hence
\[
 \mathcal U_{k+1}(u)=p\sum_{a\bmod p^k}^{*}\
 \sum_{\substack{x\bmod p^{k+1}\\F(x)\equiv0\imod p}}
 \xi_p(x)\e_{p^{k+1}}(aF(x)+u\cdot x)=p\,\mathcal S_{p,p^k}(u),
\]
and \eqref{eq:ramified} follows from \eqref{eq:Uj} with $j=k+1$.
\end{proof}

For a squarefree $q$ with $(q,N)=1$ and an integer $c\ge1$ we write
\begin{equation}\label{eq:q1q2}
 q_1=(q,c),\qquad q_2=q/q_1,
\end{equation}
as in \cite{IM}.

\begin{lemma}\label{lem:goodglobal}
Let $q\ge1$ be squarefree with $(q,N)=1$, let $c\ge1$, and let $q_1$
and $q_2$ be as in \eqref{eq:q1q2}. If $(c,N_F)=1$, then
\begin{equation}\label{eq:goodglobal}
 \mathcal S_{q,c}(u)
 =C_{q_2}(u)\,q_1c^2\chi(q_1c)\,\xi_{q_1}(\adj(A)u)\,R_{q_1c}(F^*(u)).
\end{equation}
If $F^*(u)=0$, then, for every $c\ge1$,
\begin{equation}\label{eq:coneratio}
 \mathcal S_{q,c}(u)
 =C_{q_2}(u)\,q_1^2\chi(q_1)\,\xi_{q_1}(\adj(A)u)\,S_c(u).
\end{equation}
\end{lemma}

\begin{proof}
Let $(c,N_F)=1$, and write $c=c_0\prod_{p\mid q_1}p^{k_p}$ with
$(c_0,q)=1$ and $k_p\ge1$. By Lemma~\ref{lem:crt},
\[
 \mathcal S_{q,c}(u)=C_{q_2}(u)\,S_{c_0}(u)
 \prod_{p\mid q_1}\mathcal S_{p,p^{k_p}}(u).
\]
We insert \eqref{eq:unram} for the prime powers dividing $c_0$ and
\eqref{eq:ramified} for the primes dividing $q_1$. The powers of the
primes have product $q_1c^2$, the values of $\chi$ have product
$\chi(q_1c)$, and the weights have product $\xi_{q_1}(\adj(A)u)$.
The Ramanujan sums have product $R_{q_1c}(F^*(u))$, because
$R_n(t)$ is multiplicative in $n$ and
$q_1c=c_0\prod_{p\mid q_1}p^{k_p+1}$. This proves
\eqref{eq:goodglobal}.

Let $F^*(u)=0$, and suppose first that $(c,N_F)=1$. Then
$R_{q_1c}(0)=\varphi(q_1c)=q_1\varphi(c)$, since every prime dividing
$q_1$ divides $c$, and $S_c(u)=\chi(c)c^2\varphi(c)$ by
\eqref{eq:unram}. Thus \eqref{eq:coneratio} follows from
\eqref{eq:goodglobal}. In general, write $c=c_1c_2$, where every
prime dividing $c_1$ divides $N_F$ and $(c_2,N_F)=1$. Then
$(q,c)=(q,c_2)$, and Lemma~\ref{lem:crt} gives
$\mathcal S_{q,c}(u)=S_{c_1}(u)\mathcal S_{q,c_2}(u)$ and
$S_c(u)=S_{c_1}(u)S_{c_2}(u)$, so that \eqref{eq:coneratio} for $c$
follows from the case of $c_2$.
\end{proof}

\subsection{Sums over the modulus and the singular series}\label{sec:moduli}

We now sum $\mathcal S_{q,c}(u)$ over $c$, separately for the
frequencies with $F^*(u)\ne0$ and for those with $F^*(u)=0$. In both
cases $q=q_1q_2$ is a factorization of $q$, and $c$ runs over the
integers with $(c,q)=q_1$, in accordance with \eqref{eq:q1q2}. We
write $c=c_1c_2$, where every prime dividing $c_1$ divides $N_F$ and
$(c_2,N_F)=1$. For $X\ge1$, the number of integers $c_1\le X$ of
this kind is at most $(1+\log X/\log2)^{\omega(N_F)}$.

\begin{lemma}\label{lem:genericmod}
Let $q=q_1q_2$ be squarefree with $(q,N)=1$, let $u\in\Z^4$ with
$F^*(u)\ne0$, and let $X\ge1$. Then
\begin{equation}\label{eq:genericmod}
 \sum_{\substack{c\le X\\(c,q)=q_1}}\frac{\abs{\mathcal S_{q,c}(u)}}{c^3}
 \ll_{F,\eps}\abs{C_{q_2}(u)}\,q_1^2\,\mathbf 1_{q_1\mid F^*(u)}
 \,(X\abs{F^*(u)})^\eps.
\end{equation}
\end{lemma}

\begin{proof}
Write $m=F^*(u)$. By Lemma~\ref{lem:crt} and
Lemma~\ref{lem:badprime},
\[
 \abs{\mathcal S_{q,c}(u)}=\abs{S_{c_1}(u)}\abs{\mathcal S_{q,c_2}(u)}
 \le4\abs\alpha^{1/2}c_1^3\abs{\mathcal S_{q,c_2}(u)},
\]
and $(c_2,q)=q_1$. By \eqref{eq:goodglobal} and
\eqref{eq:weightproperties},
\[
 \abs{\mathcal S_{q,c_2}(u)}\le\abs{C_{q_2}(u)}\,q_1c_2^2\abs{R_{q_1c_2}(m)}.
\]
Write $c_2=c_0\prod_{p\mid q_1}p^{k_p}$ with $(c_0,q)=1$ and
$k_p\ge1$. Since $R_n(m)$ is multiplicative in $n$,
\begin{equation}\label{eq:ramsplit}
 q_1\sum_{\substack{c_2\le X\\(c_2,q)=q_1}}\frac{\abs{R_{q_1c_2}(m)}}{c_2}
 \le\sum_{c_0\le X}\frac{\abs{R_{c_0}(m)}}{c_0}
 \prod_{p\mid q_1}\bigg(p\sum_{k\ge1}\frac{\abs{R_{p^{k+1}}(m)}}{p^k}\bigg).
\end{equation}
By \eqref{eq:ramanujan}, $R_{p^j}(m)$ is $\varphi(p^j)$ if
$p^j\mid m$, is $-p^{j-1}$ if $p^{j-1}$ is the exact power of $p$
dividing $m$, and is $0$ otherwise. Hence, with $v=v_p(m)$,
\[
 p\sum_{k\ge1}\frac{\abs{R_{p^{k+1}}(m)}}{p^k}
 =\begin{cases}0,&v=0,\\p\big((v-1)(p-1)+1\big),&v\ge1,\end{cases}
\]
which is at most $p^2v$. The product in \eqref{eq:ramsplit} therefore
vanishes unless $q_1\mid m$, and is at most $q_1^2\tau(\abs m)$. By
\eqref{eq:ramanujan} again,
\[
 \sum_{n\le X}\frac{\abs{R_n(m)}}n
 \le\sum_{h\mid m}\sum_{j\le X/h}\frac1j\le\tau(\abs m)(1+\log X).
\]
Collecting these bounds and summing over $c_1$, we find that the left
side of \eqref{eq:genericmod} is
\[
 \ll_F\abs{C_{q_2}(u)}\,q_1^2\,\mathbf 1_{q_1\mid m}\,
 \tau(\abs m)^2(1+\log X)^{1+\omega(N_F)},
\]
and the divisor bound completes the proof.
\end{proof}

When $F^*(u)=0$ we need the cancellation in the sum over $c$. It
comes from the character $\chi$, which is not principal because
$\alpha$ is not a square.

\begin{lemma}\label{lem:conemod}
Let $q=q_1q_2$ be squarefree with $(q,N_F)=1$, let $u\in\Z^4$ with
$F^*(u)=0$, and let $X\ge1$. Then
\begin{equation}\label{eq:conemod}
 \sum_{\substack{c\le X\\(c,q)=q_1}}S_c(u)\ll_{F,\eps}X^3(Xq)^\eps.
\end{equation}
\end{lemma}

\begin{proof}
Let $f$ be the conductor of $\chi$. Since $\chi$ has period $f$ and
its sum over a period vanishes, the sums of $\chi$ over intervals are
at most $f$ in absolute value. Partial summation gives, for $V>0$,
\begin{equation}\label{eq:chicube}
 \bigg|\sum_{j\le V}\chi(j)j^3\bigg|\le2fV^3.
\end{equation}
Let $M\ge1$ be an integer. Detecting the condition $(k,M)=1$ by the
M\"obius function and applying \eqref{eq:chicube}, we obtain, for
$Z>0$,
\begin{equation}\label{eq:chicoprime}
 \bigg|\sum_{\substack{k\le Z\\(k,M)=1}}\chi(k)k^3\bigg|
 =\bigg|\sum_{e\mid M}\mu(e)\chi(e)e^3\sum_{j\le Z/e}\chi(j)j^3\bigg|
 \le2^{\omega(M)+1}fZ^3.
\end{equation}

Consider first the integers $c$ with $(c,N_F)=1$. For these
$S_c(u)=\chi(c)c^2\varphi(c)$, by \eqref{eq:unram} and
Lemma~\ref{lem:crt}. The integers $c$ with $(c,q)=q_1$ are the
numbers $c=q_1n$ with $(n,q_2)=1$, and for them
\[
 \varphi(q_1n)=\varphi(q_1)\,n\sum_{\substack{h\mid n\\(h,q_1)=1}}\frac{\mu(h)}h.
\]
Writing $n=hk$, we get, for $Y\ge1$,
\[
 \sum_{\substack{c\le Y\\(c,q)=q_1\\(c,N_F)=1}}\chi(c)c^2\varphi(c)
 =\chi(q_1)q_1^2\varphi(q_1)
 \sum_{\substack{h\le Y/q_1\\(h,qN_F)=1}}\mu(h)\chi(h)h^2
 \sum_{\substack{k\le Y/(q_1h)\\(k,q_2N_F)=1}}\chi(k)k^3.
\]
By \eqref{eq:chicoprime} with $M=q_2N_F$, the inner sum is at most
$2^{\omega(q_2N_F)+1}f(Y/(q_1h))^3$ in absolute value. Hence
\begin{equation}\label{eq:conegood}
 \bigg|\sum_{\substack{c\le Y\\(c,q)=q_1\\(c,N_F)=1}}\chi(c)c^2\varphi(c)\bigg|
 \le2^{\omega(q_2N_F)+1}f\,Y^3\sum_{h\le Y}\frac1h
 \le2^{\omega(q_2N_F)+1}f\,Y^3(1+\log Y).
\end{equation}

In general write $c=c_1c_2$ as above. Then
$S_c(u)=S_{c_1}(u)S_{c_2}(u)$ and $(c,q)=(c_2,q)$. By
Lemma~\ref{lem:badprime} and \eqref{eq:conegood} with $Y=X/c_1$,
\[
 \bigg|\sum_{\substack{c\le X\\(c,q)=q_1}}S_c(u)\bigg|
 \le\sum_{c_1\le X}2^{\omega(q_2N_F)+3}\abs\alpha^{1/2}f\,c_1^3
 \Big(\frac X{c_1}\Big)^3(1+\log X).
\]
The number of $c_1$ is at most $(1+\log X/\log2)^{\omega(N_F)}$, and
$2^{\omega(q_2)}\ll_\eps q^\eps$. This proves \eqref{eq:conemod}.
\end{proof}

We turn to the singular series. Recall the local densities
$\sigma_p$ and the series $\S_F$ of \eqref{eq:singint} and
\eqref{eq:SF}.

\begin{lemma}\label{lem:density}
For every prime $p$ the limit defining $\sigma_p$ exists, and
\begin{equation}\label{eq:densityseries}
 \sigma_p=\sum_{k\ge0}\frac{S_{p^k}(0)}{p^{4k}},
\end{equation}
where the series converges absolutely. If $p\nmid N_F$, then
\begin{equation}\label{eq:sigmagood}
 \sigma_p=\frac{1-\chi(p)p^{-2}}{1-\chi(p)p^{-1}}.
\end{equation}
In particular the product in \eqref{eq:SF} converges absolutely.
\end{lemma}

\begin{proof}
Let $N_{p,k}$ be the number of $x$ modulo $p^k$ with
$F(x)\equiv0\imod{p^k}$. By orthogonality,
\[
 p^kN_{p,k}=\sum_{a\bmod p^k}\ \sum_{x\bmod p^k}\e_{p^k}(aF(x)).
\]
We sort the residues $a$ according to the power $p^{k-j}$ of $p$
dividing them, where $0\le j\le k$. The residues with a given $j$
contribute $p^{4(k-j)}S_{p^j}(0)$. Hence
\[
 p^{-3k}N_{p,k}=\sum_{j=0}^k\frac{S_{p^j}(0)}{p^{4j}}.
\]
By Lemma~\ref{lem:badprime} the terms on the right are
$O_F(p^{-j})$, so that the series in \eqref{eq:densityseries}
converges absolutely and equals $\sigma_p$. If $p\nmid N_F$, then
$S_{p^k}(0)=\chi(p)^kp^{2k}\varphi(p^k)$ by \eqref{eq:unram}, and
\[
 \sigma_p=1+\Big(1-\frac1p\Big)\sum_{k\ge1}\frac{\chi(p)^k}{p^k}
 =\frac{1-\chi(p)p^{-2}}{1-\chi(p)p^{-1}}.
\]
The factor of \eqref{eq:SF} at such a prime is $1-\chi(p)p^{-2}$,
which gives the last assertion.
\end{proof}

\begin{lemma}\label{lem:singseries}
Let $q\ge1$ be squarefree with $(q,N_F)=1$. The series
$\sum_{(c,q)=1}S_c(0)c^{-4}$ converges, and
\begin{equation}\label{eq:seriesfactor}
 \sum_{\substack{c\ge1\\(c,q)=1}}\frac{S_c(0)}{c^4}
 =\S_F\prod_{p\mid q}\sigma_p^{-1}.
\end{equation}
Moreover, for $X\ge1$,
\begin{equation}\label{eq:seriestail}
 \sum_{\substack{c>X\\(c,q)=1}}\frac{S_c(0)}{c^4}\ll_{F,\eps}q^\eps X^{-1+\eps}.
\end{equation}
\end{lemma}

\begin{proof}
Let $T(v)$ be the sum of $S_c(0)$ over the integers $c\le v$ with
$(c,q)=1$. By Lemma~\ref{lem:conemod} with $u=0$ and $q_1=1$ we have
$T(v)\ll_{F,\eps}v^{3+\eps}q^\eps$. For real $z\ge0$ and $X\ge1$,
partial summation gives
\[
 \sum_{\substack{c\le X\\(c,q)=1}}\frac{S_c(0)}{c^{4+z}}
 =\frac{T(X)}{X^{4+z}}+(4+z)\int_1^X\frac{T(v)}{v^{5+z}}\,dv.
\]
Letting $X\to\infty$, we see that the series
\[
 \mathcal D_q(z)=\sum_{\substack{c\ge1\\(c,q)=1}}\frac{S_c(0)}{c^{4+z}}
 =(4+z)\int_1^\infty\frac{T(v)}{v^{5+z}}\,dv
\]
converges for every $z\ge0$, and that it is continuous on $z\ge0$ by
dominated convergence. For $z=0$ the two formulas give
\[
 \sum_{\substack{c>X\\(c,q)=1}}\frac{S_c(0)}{c^4}
 =4\int_X^\infty\frac{T(v)}{v^5}\,dv-\frac{T(X)}{X^4},
\]
which proves \eqref{eq:seriestail}.

Let $z>0$. By Lemma~\ref{lem:badprime} the series
$\mathcal D_q(z)$ converges absolutely, and by Lemma~\ref{lem:crt}
it has the Euler product
\[
 \mathcal D_q(z)=\prod_{p\nmid q}\sigma_p(z),
 \qquad\sigma_p(z)=\sum_{k\ge0}\frac{S_{p^k}(0)}{p^{k(4+z)}}.
\]
The computation in the proof of Lemma~\ref{lem:density} gives
$\sigma_p(z)=(1-\chi(p)p^{-2-z})/(1-\chi(p)p^{-1-z})$ for
$p\nmid N_F$. Comparing with the Euler products of $L(1+z,\chi)$ and
$L(2+z,\chi)$, we obtain
\[
 \mathcal D_q(z)=\frac{L(1+z,\chi)}{L(2+z,\chi)}
 \prod_{p\mid q}\sigma_p(z)^{-1}
 \prod_{p\mid N_F}\frac{1-\chi(p)p^{-1-z}}{1-\chi(p)p^{-2-z}}\,\sigma_p(z).
\]
As $z\to0^+$, each of the finitely many local factors tends to its
value at $z=0$, where $\sigma_p(0)=\sigma_p$ by
\eqref{eq:densityseries}; for $p\mid N_F$ the series $\sigma_p(z)$
converges uniformly on $z\ge0$, by Lemma~\ref{lem:badprime}. The
quotient of $L$-functions tends to $L(1,\chi)/L(2,\chi)$: the series
$\sum_n\chi(n)n^{-s}$ converges uniformly for $s\ge1$, because the
partial sums of $\chi$ are bounded, so that $L(s,\chi)$ is continuous
there; see \cite[Section~4.3]{MV}. Since
$L(2,\chi)^{-1}=\prod_p(1-\chi(p)p^{-2})$, the limit is
\[
 L(1,\chi)\prod_p\Big(1-\frac{\chi(p)}p\Big)\sigma_p
 \prod_{p\mid q}\sigma_p^{-1}
 =\S_F\prod_{p\mid q}\sigma_p^{-1},
\]
by \eqref{eq:sigmagood} and \eqref{eq:SF}. By the continuity of
$\mathcal D_q$ at $z=0$, this proves \eqref{eq:seriesfactor}.
\end{proof}

For a squarefree $q\ge1$ with $(q,N)=1$ we set
\begin{equation}\label{eq:Sxi}
 \S_\xi(q)=\frac{C_q(0)}{q^3}\sum_{\substack{c\ge1\\(c,q)=1}}\frac{S_c(0)}{c^4}
 =\S_F\prod_{p\mid q}\frac{C_p(0)}{p^3\sigma_p},
\end{equation}
where the second expression follows from \eqref{eq:seriesfactor}.
This is the coefficient of the main term in
Section~\ref{sec:analytic}. By \eqref{eq:Cstrata} we have
$\abs{C_p(0)}\le2Kp^2$, and $\sigma_p^{-1}\le2$ by
\eqref{eq:sigmagood}. Hence
\begin{equation}\label{eq:Sxibound}
 \abs{\S_\xi(q)}\le\abs{\S_F}\frac{(4K)^{\omega(q)}}q\ll_{F,K,\eps}q^{-1+\eps}.
\end{equation}

\subsection{Averages over exceptional frequencies}\label{sec:average}

By \eqref{eq:Cstrata} the sum $C_p(u)$ is $O_K(p^{3/2})$ unless
$D_p(u)=0$, where it is $O_K(p^2)$. The two lemmas
of this section show that the larger bound is harmless on average
over $u$. For a squarefree $h\ge1$ with $(h,N)=1$ and $u\in\Z^4$ set
\[
 \Delta_h(u)=\prod_{p\mid h}\mathbf 1_{D_p(u)=0},\qquad\Delta_1(u)=1.
\]

\begin{lemma}\label{lem:box}
Let $q_1$ and $h$ be coprime squarefree positive integers with
$(q_1h,N)=1$, and let $U\ge1$. Then
\begin{equation}\label{eq:box}
 \#\{u\in\Z^4:\lVert u\rVert\le U,\ q_1\mid F^*(u),\ \Delta_h(u)=1\}
 \ll_{E,\eps}(q_1h)^\eps\Big(\frac{U^4}{q_1h}+U^3\Big),
\end{equation}
and
\begin{equation}\label{eq:weightedbox}
 \sum_{\substack{0<\lVert u\rVert\le U\\q_1\mid F^*(u)}}
 \frac{\Delta_h(u)}{\lVert u\rVert}
 \ll_{E,\eps}(q_1h)^\eps\Big(\frac{U^3}{q_1h}+U^2\Big).
\end{equation}
\end{lemma}

\begin{proof}
We apply Lemma~\ref{lem:polynomialbox} with $n=4$ and the modulus
$m=q_1h$. At a prime $p\mid q_1$ we take for $P_p$ the reduction of
$F^*$ modulo $p$. It is not the zero polynomial, because $p$ is odd
and the matrix $\adj(A)$, of determinant $\alpha^3$, is not zero
modulo $p$. At a prime $p\mid h$ we take $P_p=D_p$. The degrees are
at most $\max\{2,E\}$, and the vectors $u$ with
$\lVert u\rVert\le U$ lie in a product of four intervals containing
at most $3U$ integers each. This gives \eqref{eq:box}. For
\eqref{eq:weightedbox} we split the range of $u$ into the parts with
$2^{-j-1}U<\lVert u\rVert\le2^{-j}U$, where $j\ge0$ and $2^j\le U$,
and apply \eqref{eq:box} with $2^{-j}U$ in place of $U$ to each of
them.
\end{proof}

\begin{lemma}\label{lem:average}
Let $q=q_1q_2$ be squarefree with $(q,N)=1$, and let $U\ge1$. Then
\begin{equation}\label{eq:average}
 \sum_{\substack{0<\lVert u\rVert\le U\\q_1\mid F^*(u)}}
 \frac{\abs{C_{q_2}(u)}}{\lVert u\rVert}
 \ll_{E,K,\eps}q^\eps q_2^{3/2}\Big(\frac{U^3}{q_1}+q_2^{1/2}U^2\Big).
\end{equation}
\end{lemma}

\begin{proof}
By Lemma~\ref{lem:crt} and \eqref{eq:Cstrata},
\begin{equation}\label{eq:Cexpand}
 \abs{C_{q_2}(u)}\le\prod_{p\mid q_2}K(p^{3/2}+p^2\,\mathbf 1_{D_p(u)=0})
 =K^{\omega(q_2)}q_2^{3/2}\sum_{h\mid q_2}h^{1/2}\Delta_h(u).
\end{equation}
We insert this into the left side of \eqref{eq:average} and apply
\eqref{eq:weightedbox} to each $h$. Since
\[
 \sum_{h\mid q_2}h^{1/2}\Big(\frac{U^3}{q_1h}+U^2\Big)
 \le\tau(q_2)\Big(\frac{U^3}{q_1}+q_2^{1/2}U^2\Big),
\]
the lemma follows.
\end{proof}

\section{Fourier transforms on the quadric modulo \texorpdfstring{$p$}{p}}\label{sec:fourier}

In this section $n\ge2$, and $G\in\Z[x_1,x_2,x_3,x_4]$ is a form of
degree $n$ satisfying \eqref{eq:smooth}; the integer $N_G$ is defined
by \eqref{eq:NG}. For a prime $p\nmid N_G$ and a nonprincipal
character $\psi$ modulo $p$ such that $\psi^n$ is principal, we
consider the function $\xi_p(x)=\psi(G(x))$ on $\Fp^4$ and the sum
\begin{equation}\label{eq:Cpsi}
 C_p(u)=\sum_{\substack{x\bmod p\\F(x)\equiv0\imod p}}\psi(G(x))\e_p(u\cdot x)
 \qquad(u\in\Z^4),
\end{equation}
which is \eqref{eq:Cdef} for this function. We show that the
hypotheses \eqref{eq:weightproperties} and \eqref{eq:Cstrata} of
Section~\ref{sec:local} hold, with constants that do not depend on
the coefficients of $G$. The two cases of interest are $n=2d$ with
$\psi=\chi_p$, and $n=rd$ with $\psi=\psi_p$ of order $r$.

\subsection{Bad primes and exceptional planes}\label{sec:cover}

For a prime $p\nmid N_F$ we denote the reductions of $F$, $F^*$ and
$G$ modulo $p$ by the same letters.

\begin{lemma}\label{lem:exceptional}
Let $n\ge2$. There are numbers $C_{F,n}\ge1$ and $E_{F,n}\ge1$ with
the following properties, for every form $G\in\Z[x_1,x_2,x_3,x_4]$ of
degree $n$ satisfying \eqref{eq:smooth}.
\begin{enumerate}
\item[\textup{(i)}] The form $G$ has bad reduction at only finitely
many primes $p\nmid N_F$, and $N_G\ll_{F,n}H_G^{C_{F,n}}$.
\item[\textup{(ii)}] For every prime $p\nmid N_G$ there is a nonzero
polynomial $D_p\in\Fp[u_1,u_2,u_3,u_4]$ of total degree at most
$E_{F,n}$, with $D_p(0)=0$, such that every $u\in\Fp^4$ with
$D_p(u)\ne0$ satisfies $F^*(u)\ne0$ and
\begin{equation}\label{eq:transverse}
 \operatorname{rank}\big(\nabla F(x),\nabla G(x),u\big)=3
 \qquad\text{for every }x\in\overline{\Fp}^4\setminus\{0\}
 \text{ with }F(x)=G(x)=u\cdot x=0.
\end{equation}
\end{enumerate}
\end{lemma}

\begin{proof}
Let $\mathbf g=(g_m)$ be indeterminates indexed by the monomials
$x^m$ of degree $n$ in $x_1,\ldots,x_4$, let
$G_{\mathbf g}=\sum_mg_mx^m$ be the generic form of degree $n$, and
let $u=(u_1,\ldots,u_4)$ be further indeterminates. Consider the
following two systems of forms in $x$. The first consists of $F$,
of $G_{\mathbf g}$, and of the six minors of order $2$ of the matrix
with rows $\nabla F$ and $\nabla G_{\mathbf g}$; its coefficients lie
in $\Z[\mathbf g]$. The second consists of $F$, of $G_{\mathbf g}$,
of $u\cdot x$, and of the four minors of order $3$ of the matrix with
rows $\nabla F$, $\nabla G_{\mathbf g}$ and $u$; its coefficients lie
in $\Z[\mathbf g,u]$. All these forms have positive degree in $x$. By
the fundamental theorem of elimination theory \cite[Chapter~IX,
Theorems~3.2 and~4.1]{La}, there are polynomials
$P_1,\ldots,P_s\in\Z[\mathbf g]$ and
$R_1,\ldots,R_t\in\Z[\mathbf g,u]$ with the following property. Let
$k$ be a field, let $\gamma$ be the vector of coefficients of a form
$G_\gamma$ of degree $n$ over $k$, and let $\upsilon\in k^4$. Then
the first system, specialized at $\gamma$, has a common zero
$x\ne0$ in $\overline k^4$ if and only if $P_i(\gamma)=0$ for every
$i$; and the second system, specialized at $(\gamma,\upsilon)$, has a
common zero $x\ne0$ in $\overline k^4$ if and only if
$R_j(\gamma,\upsilon)=0$ for every $j$. These polynomials depend only
on $F$ and $n$.

(i) A common zero $x\ne0$ of the first system at $\gamma$ is a point
with $F(x)=G_\gamma(x)=0$ at which the rank of
$(\nabla F(x),\nabla G_\gamma(x))$ is less than $2$. Hence
\eqref{eq:smooth} means that $P_i(G)\ne0$ for some $i$, where
$P_i(G)$ is the value of $P_i$ at the coefficients of $G$, and bad
reduction at a prime $p\nmid N_F$ means that $p\mid P_i(G)$ for every
$i$. Fix $i$ with $P_i(G)\ne0$. The primes of bad reduction divide
the nonzero integer $P_i(G)$, so that they are finite in number, and
their product is at most $\abs{P_i(G)}\ll_{F,n}H_G^{\deg P_i}$. This
proves (i), with $C_{F,n}$ the largest of the degrees of the $P_i$.

(ii) Let $p\nmid N_G$. Then \eqref{eq:smooth} holds for the
reductions modulo $p$, over $\overline{\Fp}$. By the Jacobian
criterion, the equations $F=G=0$ define in $\PP^3$ over
$\overline{\Fp}$ a curve all of whose points are nonsingular. Its
irreducible components are therefore nonsingular projective curves.
By Bertini's theorem \cite[Chapter~II, Theorem~8.18]{Ha}, for each
component the planes that do not contain it and whose intersection
with it is regular at every point form a dense open subset of the
dual projective space. Finitely many dense open subsets have a
common point. Hence there is $u_0\in\overline{\Fp}^4\setminus\{0\}$
such that the plane $u_0\cdot x=0$ has this property for every
component. The tangent line to the curve at a point $x$ is defined by
$\nabla F(x)\cdot y=\nabla G(x)\cdot y=0$, and the intersection with
the plane is regular at $x$ exactly when the plane does not contain
this line, that is, when $u_0$ is not a linear combination of
$\nabla F(x)$ and $\nabla G(x)$. Thus the second system, specialized
at the coefficients of $G$ and at $u_0$, has no common zero
$x\ne0$, and $R_j(G,u_0)\ne0$ for some $j$. For this $j$, the
polynomial $R_j(G,u)\in\Fp[u_1,\ldots,u_4]$, obtained by reducing
the coefficients of $G$ modulo $p$, is not zero. Set
\[
 D_p(u)=F^*(u)\,R_j(G,u).
\]
Since $p\nmid N_F$, the form $F^*$ is not zero modulo $p$, so that
$D_p\ne0$. Its degree is at most $2$ plus the largest of the degrees
of the $R_j$ in $u$, which we call $E_{F,n}$, and $D_p(0)=0$. If
$u\in\Fp^4$ and $D_p(u)\ne0$, then $F^*(u)\ne0$ and $R_j(G,u)\ne0$,
so that the second system at $u$ has no common zero $x\ne0$ in
$\overline{\Fp}^4$. This is \eqref{eq:transverse}.
\end{proof}

With $n=2d$, part (i) is the first assertion of
Theorem~\ref{thm:polynomial}, and with $n=rd$ it is the corresponding
assertion of Theorem~\ref{thm:cyclic}.

\subsection{Plane sections}\label{sec:hyperplane}

For $u\in\Fp^4$, not zero, the points of $\PP^3$ with
$F(x)=u\cdot x=0$ form the section of the quadric by the plane
$u\cdot x=0$. When $F^*(u)\ne0$ it is a conic that can be
parametrized by the projective line.

\begin{lemma}\label{lem:conic}
Let $p\nmid N_F$ be a prime and let $u\in\Fp^4$ with $F^*(u)\ne0$.
There is a vector $\mu=(\mu_1,\mu_2,\mu_3,\mu_4)$ of binary quadratic
forms $\mu_i\in\Fp[s,t]$ with the following properties.
\begin{enumerate}
\item[\textup{(i)}] For every field $k$ containing $\Fp$, the map
$(s:t)\mapsto\mu(s,t)$ is a bijection from $\PP^1(k)$ onto the set of
points $x$ of $\PP^3(k)$ with $F(x)=u\cdot x=0$.
\item[\textup{(ii)}] For every $(s,t)\in\overline{\Fp}^2$ other than
$(0,0)$, the vectors $\partial\mu/\partial s$ and
$\partial\mu/\partial t$ at $(s,t)$ are linearly independent.
\end{enumerate}
\end{lemma}

\begin{proof}
Let $V$ be the space of $x\in\Fp^4$ with $u\cdot x=0$, and write
$b(x,y)=x^{\mathsf T}Ay$. Suppose that $x\in V$ and $b(x,y)=0$ for
every $y\in V$. Then $Ax$ is a multiple of $u$, so that
$x=c\adj(A)u$ for some $c\in\Fp$, and $u\cdot x=cF^*(u)$. Hence
$c=0$. Thus the restriction of $F$ to $V$ is a nondegenerate
quadratic form in three variables. In suitable coordinates on $V$ it
is $a_1y_1^2+a_2y_2^2+a_3y_3^2$ with $a_1a_2a_3\ne0$. The sets of
values of $a_1y_1^2$ and of $-a_3-a_2y_2^2$ have $(p+1)/2$ elements
each, so that they meet, and $F$ has a zero $e_1\ne0$ in $V$. As in
the proof of Lemma~\ref{lem:quadriccount}, there are $e_2,e_3\in V$
such that $e_1,e_2,e_3$ is a basis of $V$ and
\[
 F(y_1e_1+y_2e_2+y_3e_3)=y_1y_2+ay_3^2,
\]
where $a\ne0$ because the form is nondegenerate. Set
\[
 \mu(s,t)=-as^2e_1+t^2e_2+ste_3.
\]
Let $k$ be a field containing $\Fp$. The points of $\PP^3(k)$ with
$u\cdot x=0$ are the points $y_1e_1+y_2e_2+y_3e_3$ with
$(y_1:y_2:y_3)\in\PP^2(k)$. If $y_1y_2+ay_3^2=0$ and $y_2\ne0$, we
may take $y_2=1$, and the point is $\mu(y_3,1)$. If $y_2=0$, then
$y_3=0$, and the point is $\mu(1,0)$. Conversely $F(\mu(s,t))=0$, and
the points $\mu(s,1)$ and $\mu(1,0)$ are distinct. This proves (i).
In the basis $e_1,e_2,e_3$ the two derivatives have the coordinates
$(-2as,0,t)$ and $(0,2t,s)$. They are linearly independent if
$s\ne0$, and also if $s=0$ and $t\ne0$, because $p$ is odd.
\end{proof}

\begin{lemma}\label{lem:section}
Let $p\nmid N_F$ be a prime and $n\ge2$. Let $\psi$ be a nonprincipal
character modulo $p$ such that $\psi^n$ is principal, let
$G\in\Fp[x_1,x_2,x_3,x_4]$ be a form of degree $n$, and let
$u\in\Fp^4$ with $F^*(u)\ne0$ satisfy \eqref{eq:transverse}. Then
\begin{equation}\label{eq:section}
 \bigg|\sum_{\substack{P\in\Qc(\Fp)\\u\cdot P=0}}\psi(G(P))\bigg|
 \le(2n-1)\sqrt p+1.
\end{equation}
\end{lemma}

\begin{proof}
Let $\mu$ be as in Lemma~\ref{lem:conic}, and set
$g(s,t)=G(\mu(s,t))$, a binary form of degree $2n$ over $\Fp$. By
Lemma~\ref{lem:conic}(i) the sum in \eqref{eq:section} is the sum of
$\psi(g(s,t))$ over $(s:t)\in\PP^1(\Fp)$. We claim that $g$ has $2n$
distinct zeros in $\PP^1(\overline{\Fp})$; the bound is then
Corollary~\ref{cor:weilline} with $2n$ in place of $n$.

Over $\overline{\Fp}$ a nonzero binary form is a product of linear
forms. If the claim is false, then $g=0$ or $g$ is divisible by the
square of a linear form, and in both cases there is
$(s_0,t_0)\ne(0,0)$ in $\overline{\Fp}^2$ at which $g$ and its two
partial derivatives vanish. Let $x_0=\mu(s_0,t_0)$, and let $y_1$ and
$y_2$ be the derivatives of $\mu$ with respect to $s$ and to $t$ at
$(s_0,t_0)$. Then $x_0\ne0$ and $F(x_0)=G(x_0)=u\cdot x_0=0$.
Differentiating the identities $F(\mu(s,t))=0$ and
$u\cdot\mu(s,t)=0$, and the form $g$, we get
\[
 \nabla F(x_0)\cdot y_i=0,\qquad u\cdot y_i=0,\qquad
 \nabla G(x_0)\cdot y_i=0\qquad(i=1,2).
\]
By Lemma~\ref{lem:conic}(ii) the vectors $y_1$ and $y_2$ are linearly
independent, so that the vectors orthogonal to both form a space of
dimension $2$. It contains $\nabla F(x_0)$, $\nabla G(x_0)$ and $u$,
which contradicts \eqref{eq:transverse}.
\end{proof}

We can now verify the hypotheses of Section~\ref{sec:local}.

\begin{lemma}\label{lem:fourier}
Let $n\ge2$, let $G$ be as at the beginning of this section, and for
every prime $p\nmid N_G$ let $D_p$ be the polynomial of
Lemma~\ref{lem:exceptional}. Let $p\nmid N_G$, and let $\psi$ be a
nonprincipal character modulo $p$ such that $\psi^n$ is principal.
Then the function $\xi_p(x)=\psi(G(x))$ satisfies
\eqref{eq:weightproperties}, and the sum \eqref{eq:Cpsi} satisfies
\begin{equation}\label{eq:fourier}
 \abs{C_p(u)}\le2n^2(p^{3/2}+p^2\,\mathbf 1_{D_p(u)=0})
 \quad(u\in\Z^4),\qquad
 C_p(0)=(p-1)\sum_{P\in\Qc(\Fp)}\psi(G(P)).
\end{equation}
\end{lemma}

\begin{proof}
We have $\psi(G(0))=\psi(0)=0$ and $\abs{\psi}\le1$, and
$\psi(G(tx))=\psi(t^n)\psi(G(x))=\psi(G(x))$ for $t\ne0$. This is
\eqref{eq:weightproperties}. Let $a$ be the sum of $\psi(G(P))$ over
$P\in\Qc(\Fp)$ and, for $u$ not divisible by $p$, let $a_u$ be the
sum over the points $P$ with $u\cdot P=0$. In \eqref{eq:Cpsi} the
vector $x=0$ contributes nothing, and the other $x$ are the multiples
$tx_P$ of a representative $x_P$ of a point $P\in\Qc(\Fp)$, with
$t\in\Fp\setminus\{0\}$. The sum of $\e_p(tu\cdot x_P)$ over $t$ is
$p-1$ if $u\cdot x_P=0$, and is $-1$ otherwise. Hence
\begin{equation}\label{eq:rays}
 C_p(u)=pa_u-a\quad(p\nmid u),\qquad C_p(u)=(p-1)a\quad(p\mid u).
\end{equation}
Since $p\nmid N_G$, Lemma~\ref{lem:katz} gives
$\abs a\le(2n^2-4n+4)p$. If $p\mid u$, then $D_p(u)=0$ and
$\abs{C_p(u)}\le(2n^2-4n+4)p^2$.

Let $D_p(u)\ne0$. By Lemma~\ref{lem:exceptional}(ii) and
Lemma~\ref{lem:section} we have $\abs{a_u}\le(2n-1)\sqrt p+1$, and
\[
 \abs{C_p(u)}\le(2n-1)p^{3/2}+(2n^2-4n+5)p\le(2n^2-2n+4)p^{3/2}.
\]

Let $D_p(u)=0$ and $p\nmid u$. The form $F$ does not vanish
identically on the space $V$ of $x\in\Fp^4$ with $u\cdot x=0$:
otherwise $2x^{\mathsf T}Ay=F(x+y)-F(x)-F(y)=0$ for all $x,y\in V$,
so that the space
$AV$, of dimension $3$, would consist of multiples of $u$. Let $P_0$
be a point of $\PP^3(\Fp)$ with $u\cdot P_0=0$ and $F(P_0)\ne0$. Each
of the $p+1$ lines through $P_0$ in the plane $u\cdot x=0$ contains
at most two zeros of $F$, because the restriction of $F$ to it is a
binary quadratic form that is not zero. Hence
$\abs{a_u}\le2p+2$, and
\[
 \abs{C_p(u)}\le2p^2+(2n^2-4n+6)p\le(2n^2-4n+8)p^2.
\]
Since $n\ge2$, the three bounds are at most $2n^2p^{3/2}$ or
$2n^2p^2$, as required.
\end{proof}

Thus the hypotheses of Section~\ref{sec:local} hold for the functions
$\xi_p(x)=\psi(G(x))$, with $N=N_G$, with $E=E_{F,n}$ and with
$K=2n^2$.

\subsection{The diagonal case}\label{sec:diagonal}

In this subsection $F$ and $G$ are the diagonal quadratic forms of
\eqref{eq:diagonal}, with $\kappa_i$ and $M$ as in
\eqref{eq:kappadef}, and $\psi=\chi_p$. We prove the first two
assertions of Theorem~\ref{thm:diagonal} and evaluate $C_p(u)$ for
every $u$.

\begin{lemma}\label{lem:diagbad}
The form $G$ satisfies \eqref{eq:smooth}, and it has bad reduction at
a prime $p\nmid N_F$ if and only if $p\mid M$. Hence
$N_G=\rad(2\alpha M)$.
\end{lemma}

\begin{proof}
Let $k$ be $\overline\Q$, or $\overline{\Fp}$ for a prime
$p\nmid N_F$. The minors of order $2$ of the matrix with rows
$\nabla F(x)$ and $\nabla G(x)$ are the products
$4x_ix_j(a_ib_j-a_jb_i)$ with $i<j$. Suppose that $M$ is not zero in
$k$, and let $x\in k^4$ be a point at which all these minors vanish.
Then at most one coordinate $x_i$ is not zero, and $F(x)=a_ix_i^2$.
Hence $F(x)=0$ only for $x=0$. This proves \eqref{eq:smooth}, and
that $G$ has good reduction at $p$ if $p\nmid M$. If $p\mid M$, then
$a_ib_j\equiv a_jb_i\imod p$ for some $i<j$. Let
$x\in\overline{\Fp}^4$ have $x_i=1$ and
$x_j^2=-a_i\overline{a_j}$, the other coordinates being zero. Then
$F(x)=0$ and $G(x)=\overline{a_j}(a_jb_i-a_ib_j)=0$, and all the
minors vanish at $x$. Thus $G$ has bad reduction at $p$.
\end{proof}

Let $p\nmid2\alpha M$ be a prime. For $1\le i\le4$ let $t_i$ be the
residue class modulo $p$ with $a_it_i+b_i\equiv0\imod p$; the four
classes are distinct because $p\nmid M$, so that $p\ge5$. For
$u\in\Z^4$ and an integer $t$ set
\begin{equation}\label{eq:pencilforms}
 \Phi_u(t)=\sum_{i=1}^4u_i^2\prod_{j\ne i}(b_j+ta_j),
 \qquad
 \nu_i(u)=\sum_{j\ne i}u_j^2\prod_{l\ne i,j}(a_ib_l-a_lb_i).
\end{equation}
Thus $\Phi_u(t)$ is the value at $u$ of the dual form of $G+tF$; it
is a polynomial in $t$ of degree at most $3$, and the coefficient of
$t^3$ is $F^*(u)$.

\begin{lemma}\label{lem:pencil}
Let $p\nmid2\alpha M$ be a prime and $u\in\Z^4$. Then
\begin{equation}\label{eq:exactpencil}
 C_p(u)=p\sum_{\substack{t\bmod p\\t\not\equiv t_1,\ldots,t_4}}
 \chi_p\big(-\Phi_u(t)\big)
 +p\sum_{\substack{1\le i\le4\\p\mid u_i}}\chi_p(\kappa_i)R_p(\nu_i(u)).
\end{equation}
In particular \eqref{eq:diagtrace} holds, and
\begin{equation}\label{eq:diagbound}
 \abs{C_p(u)}\le2p^{3/2}+4p\qquad\text{if }p\nmid F^*(u)u_1u_2u_3u_4.
\end{equation}
\end{lemma}

\begin{proof}
For $x\in\Z^4$ the sum of $\e_p(vF(x))$ over $v$ modulo $p$ is $p$
if $p\mid F(x)$, and is $0$ otherwise. With $\tau_p$ as in the proof
of Lemma~\ref{lem:gauss}, we have
\begin{equation}\label{eq:gaussinv}
 \sum_{s\bmod p}\chi_p(s)\e_p(sm)=\chi_p(m)\tau_p
\end{equation}
for every integer $m$: both sides vanish if $p\mid m$, and otherwise
we replace $s$ by $s\overline m$. We apply this with $m=G(x)$, and
write $v=st$ for $p\nmid s$. Then \eqref{eq:Cpsi} becomes
\begin{equation}\label{eq:pencilsum}
 C_p(u)=\frac1{p\tau_p}\sum_{s\bmod p}\chi_p(s)\sum_{t\bmod p}\
 \prod_{i=1}^4\ \sum_{y\bmod p}\e_p\big(s(b_i+ta_i)y^2+u_iy\big).
\end{equation}
Let $c$ and $w$ be integers. If $p\nmid c$, completing the square and
the evaluation of the Gauss sum in the proof of
Lemma~\ref{lem:gauss} give
\[
 \sum_{y\bmod p}\e_p(cy^2+wy)=\chi_p(c)\tau_p\,\e_p(-\overline{4c}\,w^2).
\]
If $p\mid c$, the sum is $p$ when $p\mid w$, and is $0$ otherwise.

Let $t\not\equiv t_i$ for every $i$, and let $\Delta$ be the product
of the four numbers $b_i+ta_i$. The product over $i$ in
\eqref{eq:pencilsum} is
$\tau_p^4\chi_p(\Delta)\e_p(-\overline{4s}\,\Phi_u(t)\overline\Delta)$.
Replacing $s$ by $\overline s$ and using \eqref{eq:gaussinv}, we see
that its sum over $s$ against $\chi_p(s)$ is
$\tau_p^5\chi_p(\Delta)\chi_p(-\Phi_u(t)\Delta)$. Since
$\tau_p^4=p^2$, the contribution of $t$ to $C_p(u)$ is
$p\chi_p(-\Phi_u(t))$.

Let $t\equiv t_i$. The factor of index $i$ is $p$ if $p\mid u_i$, and
is $0$ otherwise. Let $\Delta_i$ be the product of the three numbers
$b_j+t_ia_j$ with $j\ne i$, and let $w_i$ be the sum of
$u_j^2\overline{(b_j+t_ia_j)}$ over $j\ne i$. The three other factors
have the product
$\tau_p^3\chi_p(s)\chi_p(\Delta_i)\e_p(-\overline{4s}\,w_i)$, and the
sum of $\e_p(-\overline{4s}\,w_i)$ over $s\not\equiv0$ is $R_p(w_i)$.
Since $\tau_p^2=\chi_p(-1)p$, by the proof of
Lemma~\ref{lem:gauss}, the contribution of $t_i$ is
$p\chi_p(-\Delta_i)R_p(w_i)$ when $p\mid u_i$. Now
$a_i(b_j+t_ia_j)\equiv a_ib_j-a_jb_i\imod p$. Hence
$-\Delta_i\equiv\kappa_i\overline{a_i}^4$, and $w_i$ is the product
of $\nu_i(u)$ and an integer prime to $p$. This proves
\eqref{eq:exactpencil}.

For $u=0$ the first sum in \eqref{eq:exactpencil} vanishes and
$R_p(\nu_i(0))=p-1$, so that $C_p(0)$ is $p(p-1)$ times the sum of
the $\chi_p(\kappa_i)$. Since $p\nmid N_G$ by
Lemma~\ref{lem:diagbad}, this and \eqref{eq:fourier} give
\eqref{eq:diagtrace}.

Let $p\nmid F^*(u)u_1u_2u_3u_4$. The second sum in
\eqref{eq:exactpencil} is empty, and $\Phi_u$ is a cubic polynomial
modulo $p$. If it has three distinct zeros in $\overline{\Fp}$, then
the sum of $\chi_p(-\Phi_u(t))$ over all $t$ modulo $p$ is at most
$2\sqrt p$ in absolute value, by Lemma~\ref{lem:weil}. Otherwise
$\Phi_u$ has a repeated zero $z$, which lies in $\Fp$ because $z$ is
a zero of the greatest common divisor of $\Phi_u$ and its
derivative, a polynomial over $\Fp$ equal to a constant multiple of
$t-z$ or of $(t-z)^2$, and $p\ge5$. Then $\Phi_u(t)$ is
$F^*(u)(t-z)^3$ or $F^*(u)(t-z)^2(t-z')$ with $z'\in\Fp$ and
$z'\ne z$, and the sum over all $t$ is $0$ or
$-\chi_p(-F^*(u)(z-z'))$. In every case the sum over
$t\not\equiv t_1,\ldots,t_4$ is at most $2\sqrt p+4$ in absolute
value, which is \eqref{eq:diagbound}.
\end{proof}

Lemma~\ref{lem:diagbad} and \eqref{eq:diagtrace} are the first two
assertions of Theorem~\ref{thm:diagonal}. By \eqref{eq:fourier} with
$n=2$ and by \eqref{eq:diagbound}, in the diagonal case the hypothesis
\eqref{eq:Cstrata} holds with $K=8$ and with the polynomial
$F^*(u)u_1u_2u_3u_4$ in place of $D_p$.

\begin{remark}\label{rem:literature}
The formula \eqref{eq:exactpencil} shows that $C_p(u)$ can have order
$p^2$. At $u=0$ it is $p(p-1)$ times the sum of the
$\chi_p(\kappa_i)$, which is $\pm2$ when $\chi(p)=-1$. For an example
with $p\nmid F^*(u)$, let
\[
 F=x_1^2+x_2^2+x_3^2-2x_4^2,\qquad
 G=x_1^2+2x_2^2+3x_3^2+5x_4^2,\qquad u=(1,0,0,3).
\]
Here $\alpha=-2$, $M=1386$ and $F^*(u)=7$, and
$\Phi_u(t)=7(t+2)^2(t+3)$, $\nu_2(u)=0$, $\nu_3(u)=7$. For every
prime $p>11$, \eqref{eq:exactpencil} gives
\[
 C_p(u)=p\big(p-1-\chi_p(-7)-\chi_p(-14)-\chi_p(-22)-\chi_p(-154)\big)
 \ge p(p-5).
\]
In the notation of \cite{M}, $C_p(u)$ is the sum $S_{p,1}(u)$, for
which \cite[Lemma~5.3]{M} states the bound $O(p^{3/2})$ for every
$u$. The sum $S_{p,p}(u)$ of \cite{IM} is
\[
 pC_p(u)-\sum_{x\bmod p}\chi_p(G(x))\e_p(u\cdot x),
\]
and the last sum, which is $S_{p,1}(u)$ in \cite{IM}, has absolute
value at most $p^2$ by \cite[Lemma~4.3]{IM}. In the example the sum
$S_{p,p}(u)$ of \cite{IM} therefore has absolute value at least
$p^2(p-6)$, while \cite[Lemma~4.6]{IM} states the bound
$O(p^{5/2})$ for every $u$. For the same forms and $u=(1,1,0,2)$,
where $F^*(u)=0$, Lemma~\ref{lem:ramified} gives
$\mathcal S_{p,p}(u)=\chi_p(2)p^4(p-1)$, while
\cite[Lemma~5.4]{M} states the bound $O(p^{9/2})$ for this sum. The
nonzero frequencies at which $C_p(u)$ is large are treated on
average, in Lemma~\ref{lem:average}.
\end{remark}

\section{Proofs of the main results}\label{sec:proofs}

\subsection{The analytic proposition}\label{sec:analytic}

In this subsection the notation and the hypotheses of
Section~\ref{sec:local} are in force. For a weight $w$, a squarefree
integer $q\ge1$ with $(q,N)=1$, and $B>0$ set
\begin{equation}\label{eq:Tw}
 T_q(B;w)=\sum_{\substack{x\in\Z^4\\F(x)=0}}w(x/B)\,\xi_q(x).
\end{equation}

\begin{proposition}\label{prop:analytic}
Let $W$ be a weight. For every squarefree $q\ge1$ with $(q,N)=1$ and
every $B\ge2$,
\begin{equation}\label{eq:analytic}
 T_q(B;W)=B^2\sigma_\infty(W,F)\,\S_\xi(q)
 +O_{F,E,K,W,\eps}\big((Bq)^{1+\eps}\big).
\end{equation}
In particular the implied constant does not depend on $N$, on the
polynomials $D_p$ or on the functions $\xi_p$.
\end{proposition}

The proof follows the method of \cite{HB} and occupies the rest of
this subsection. We begin with the transformation of $T_q(B;w)$ by
the delta symbol and the Poisson summation formula.

\begin{lemma}\label{lem:poisson}
Let $w$ be a weight, let $q\ge1$ be squarefree with $(q,N)=1$, and
let $B>0$ be such that $Q=B/\sqrt q\ge3$. Then
\begin{equation}\label{eq:poisson}
 T_q(B;w)=c_Q\frac{B^2}{q^3}\sum_{c\ge1}\frac1{c^4}\sum_{u\in\Z^4}
 \mathcal S_{q,c}(u)\,I_{c/Q}^*\Big(\frac u{\sqrt q};w\Big),
\end{equation}
where only finitely many $c$ contribute, and the sum over $u$
converges absolutely for each of them.
\end{lemma}

\begin{proof}
Let $x\in\Z^4$ with $F(x)\equiv0\imod q$. We apply \eqref{eq:delta}
to the integer $n=F(x)/q$. Since $n/Q^2=F(x/B)$ and
$\e_c(an)=\e_{qc}(aF(x))$, this gives
\[
 T_q(B;w)=\frac{c_Q}{Q^2}\sum_{c\ge1}\ \sum_{a\bmod c}^{*}\
 \sum_{\substack{x\in\Z^4\\F(x)\equiv0\imod q}}
 \xi_q(x)\,\e_{qc}(aF(x))\,g_c(x),
 \qquad g_c(x)=w\Big(\frac xB\Big)h\bigg(\frac cQ,F\Big(\frac xB\Big)\bigg).
\]
The function $g_c$ is infinitely differentiable with compact support,
and by the last assertion of Lemma~\ref{lem:delta} it vanishes
identically when $c/Q$ exceeds $1+2\max\abs{F(y)}$, the maximum being
over the support of $w$. We sort the vectors $x$ according to their
residue classes $z$ modulo $qc$. The numbers $\xi_q(x)$ and
$\e_{qc}(aF(x))$ and the condition $F(x)\equiv0\imod q$ depend only
on $z$. By the Poisson summation formula in $\R^4$, as in the proof
of \cite[Theorem~2]{HB},
\[
 \sum_{m\in\Z^4}g_c(z+qcm)=\frac1{(qc)^4}\sum_{u\in\Z^4}
 \e_{qc}(u\cdot z)\int_{\R^4}g_c(x)\,\e_{qc}(-u\cdot x)\,dx,
\]
where the sum over $u$ converges absolutely. The substitution $x=By$
shows that the integral is
\[
 B^4\int_{\R^4}w(y)\,h\Big(\frac cQ,F(y)\Big)
 \e\Big(-\frac{Bu\cdot y}{qc}\Big)\,dy
 =B^4I_{c/Q}^*\Big(\frac u{\sqrt q};w\Big),
\]
by \eqref{eq:Istar}, because $B/(qc)=(Q/c)q^{-1/2}$. Summing over $z$
and $a$, we obtain the sum $\mathcal S_{q,c}(u)$ of
\eqref{eq:Snewdef}, and \eqref{eq:poisson} follows because
$B^4Q^{-2}q^{-4}=B^2q^{-3}$.
\end{proof}

Let $\mathcal W$ be a set of weights as in Lemma~\ref{lem:integrals},
and let $\delta$ and $r_0$ be the numbers given by that lemma. In
Lemmas~\ref{lem:offcone}, \ref{lem:oncone} and~\ref{lem:zerofreq}
below, $w=W_y$ is one of the functions \eqref{eq:pieces}, with
$W\in\mathcal W$ and $y\in\R^4$, such that $F$ has a zero in the
support of $w$; thus \eqref{eq:Izero}, \eqref{eq:Irapid} and
\eqref{eq:Idecay} hold for $w$. Further $0<\eta\le1/4$, the integer
$q\ge1$ is squarefree with $(q,N)=1$, and $B\ge9$ and $q\le B$. We
write
\[
 Q=\frac B{\sqrt q},\qquad U=\sqrt q\,B^\eta,\qquad
 I_c(u)=I_{c/Q}^*\Big(\frac u{\sqrt q};w\Big),
\]
so that $3\le\sqrt B\le Q\le B$ and $U\ge1$. In these three lemmas
the implied constants depend at most on $F$, $\mathcal W$, $E$, $K$
and $\eta$. Since $I_r^*(v;w)=0$ for $r\ge r_0$,
Lemma~\ref{lem:poisson} gives
\begin{equation}\label{eq:poissonpiece}
 T_q(B;w)=c_Q\frac{B^2}{q^3}\sum_{c<r_0Q}\frac1{c^4}\sum_{u\in\Z^4}
 \mathcal S_{q,c}(u)\,I_c(u).
\end{equation}
Here $c_Q\ll1$, by Lemma~\ref{lem:delta}. For fixed $v$ the integral
$I_r^*(v;w)$ is a continuously differentiable function of $r>0$,
because $h$ is infinitely differentiable and $w$ has compact support.
Let $t\ge1$ be real and $0<\lVert u\rVert\le U$. With $r=t/Q$ and
$v=u/\sqrt q$, the bound \eqref{eq:Idecay} with $\eps=\eta$ gives
\begin{equation}\label{eq:Ipiece}
 \abs{I_r^*(v;w)}+r\,\bigg|\frac{\partial I_r^*(v;w)}{\partial r}\bigg|
 \ll\Big(\frac tQ\Big)^{1-2\eta}
 \bigg(\frac{\sqrt q}{\lVert u\rVert}\bigg)^{1-\eta}
 \le\frac{tq}{B\lVert u\rVert}\,B^{3\eta},
\end{equation}
because $(Q/t)^{2\eta}\le B^{2\eta}$,
$(\lVert u\rVert/\sqrt q)^\eta\le B^\eta$ and $\sqrt q/Q=q/B$.

\begin{lemma}\label{lem:offcone}
The terms of \eqref{eq:poissonpiece} with $\lVert u\rVert>U$
contribute $O(1)$, and those with $0<\lVert u\rVert\le U$ and
$F^*(u)\ne0$ contribute $O((Bq)^{1+8\eta})$.
\end{lemma}

\begin{proof}
By \eqref{eq:Snewdef} and \eqref{eq:weightproperties} we have
$\abs{\mathcal S_{q,c}(u)}\le q^4c^5$, and by \eqref{eq:Irapid},
\[
 I_c(u)\ll_j\frac Qc\bigg(\frac{\sqrt q}{\lVert u\rVert}\bigg)^j
\]
for $u\ne0$ and every integer $j\ge1$. For $j\ge5$ the sum of
$\lVert u\rVert^{-j}$ over the $u\in\Z^4$ with $\lVert u\rVert>U$ is
$O_j(U^{4-j})$, and $q^{j/2}U^{4-j}=q^2B^{-\eta(j-4)}$. Hence the
terms with $\lVert u\rVert>U$ contribute
\[
 \ll_j\frac{B^2}{q^3}\sum_{c<r_0Q}q^4Q\,q^2B^{-\eta(j-4)}
 \ll B^2q^3Q^2B^{-\eta(j-4)}=B^4q^2B^{-\eta(j-4)}.
\]
Since $q\le B$, this is $O(1)$ for $j\ge4+6/\eta$.

Let $0<\lVert u\rVert\le U$. By \eqref{eq:Ipiece} with $t=c$, the
terms of \eqref{eq:poissonpiece} with these $u$ and with
$F^*(u)\ne0$ contribute
\[
 \ll\frac B{q^2}B^{3\eta}\sum_{q_1q_2=q}\
 \sum_{\substack{0<\lVert u\rVert\le U\\F^*(u)\ne0}}\frac1{\lVert u\rVert}
 \sum_{\substack{c\le r_0Q\\(c,q)=q_1}}\frac{\abs{\mathcal S_{q,c}(u)}}{c^3},
\]
where we have sorted the integers $c$ according to the value $q_1$
of $(c,q)$. We apply Lemma~\ref{lem:genericmod} with $\eps=\eta$ and
$X=r_0Q$. Since $\abs{F^*(u)}\ll_FU^2$ and
$QU^2\le BqB^{2\eta}$, the factor $(X\abs{F^*(u)})^\eta$ is
$O((Bq)^{2\eta})$. The contribution is therefore
\[
 \ll\frac B{q^2}B^{3\eta}(Bq)^{2\eta}\sum_{q_1q_2=q}q_1^2
 \sum_{\substack{0<\lVert u\rVert\le U\\q_1\mid F^*(u)}}
 \frac{\abs{C_{q_2}(u)}}{\lVert u\rVert}.
\]
By Lemma~\ref{lem:average} with $\eps=\eta$, and since
$U^3=q^{3/2}B^{3\eta}$ and $U^2=qB^{2\eta}$,
\[
 q_1^2\sum_{\substack{0<\lVert u\rVert\le U\\q_1\mid F^*(u)}}
 \frac{\abs{C_{q_2}(u)}}{\lVert u\rVert}
 \ll q^\eta(q_1q_2^{3/2}U^3+q_1^2q_2^2U^2)
 \le2q^{3+\eta}B^{3\eta}.
\]
There are $\tau(q)\ll q^\eta$ factorizations $q=q_1q_2$. Altogether
these terms contribute
\[
 \ll\frac B{q^2}B^{3\eta}(Bq)^{2\eta}q^{3+2\eta}B^{3\eta}
 \le(Bq)^{1+8\eta}.\qedhere
\]
\end{proof}

On the cone $F^*(u)=0$ the sum over $c$ is carried out before
absolute values are taken.

\begin{lemma}\label{lem:oncone}
The terms of \eqref{eq:poissonpiece} with $0<\lVert u\rVert\le U$
and $F^*(u)=0$ contribute $O(B\sqrt q\,(Bq)^{6\eta})$.
\end{lemma}

\begin{proof}
Let $u$ be such a vector, and let $q=q_1q_2$. By
\eqref{eq:coneratio}, the terms of \eqref{eq:poissonpiece} with this
$u$ and with $(c,q)=q_1$ have the sum
\begin{equation}\label{eq:conesum}
 c_Q\frac{B^2}{q^3}\,C_{q_2}(u)\,q_1^2\chi(q_1)\,\xi_{q_1}(\adj(A)u)
 \sum_{\substack{c<r_0Q\\(c,q)=q_1}}S_c(u)\,\phi(c),
 \qquad\phi(t)=t^{-4}I_{t/Q}^*\Big(\frac u{\sqrt q};w\Big).
\end{equation}
Let $T_u(t)$ be the sum of $S_c(u)$ over the integers $c\le t$ with
$(c,q)=q_1$. Since $\phi(t)=0$ for $t\ge r_0Q$, partial summation
gives
\[
 \sum_{\substack{c<r_0Q\\(c,q)=q_1}}S_c(u)\,\phi(c)
 =-\int_1^{r_0Q}T_u(t)\,\phi'(t)\,dt.
\]
Here $T_u(t)\ll t^3(tq)^\eta$ by Lemma~\ref{lem:conemod}. With
$r=t/Q$ and $v=u/\sqrt q$ we have
\[
 \phi'(t)=-4t^{-5}I_r^*(v;w)
 +t^{-5}\,r\,\frac{\partial I_r^*(v;w)}{\partial r},
\]
so that $\phi'(t)\ll t^{-4}q(B\lVert u\rVert)^{-1}B^{3\eta}$ by
\eqref{eq:Ipiece}. Hence
\[
 \sum_{\substack{c<r_0Q\\(c,q)=q_1}}S_c(u)\,\phi(c)
 \ll\frac q{B\lVert u\rVert}B^{3\eta}q^\eta\int_1^{r_0Q}t^{-1+\eta}\,dt
 \ll\frac q{B\lVert u\rVert}B^{3\eta}(Bq)^\eta.
\]
By \eqref{eq:Cstrata} and Lemma~\ref{lem:crt} we have
$\abs{C_{q_2}(u)}\le(2K)^{\omega(q_2)}q_2^2\ll q_2^2q^\eta$. The
expression \eqref{eq:conesum} is therefore
\[
 \ll\frac{B^2}{q^3}\,q_2^2q^\eta\,q_1^2\,
 \frac q{B\lVert u\rVert}B^{3\eta}(Bq)^\eta
 =\frac B{\lVert u\rVert}B^{3\eta}(Bq)^\eta q^\eta.
\]
We sum this over the $\tau(q)\ll q^\eta$ factorizations of $q$ and
over $u$. By \eqref{eq:conecount} with $\eps=\eta$, the sum of
$\lVert u\rVert^{-1}$ over the $u$ in question is
$O(U^{1+\eta})$, and $U^{1+\eta}\le\sqrt q\,q^\eta B^{2\eta}$. The
total is $O(B\sqrt q\,(Bq)^{6\eta})$.
\end{proof}

\begin{lemma}\label{lem:zerofreq}
The terms of \eqref{eq:poissonpiece} with $u=0$ contribute
\[
 B^2\sigma_\infty(w,F)\,\S_\xi(q)+O\big(B(Bq)^{2\eta}\big).
\]
\end{lemma}

\begin{proof}
Let $q_1=(c,q)$ and $q_2=q/q_1$. By \eqref{eq:coneratio} with $u=0$,
\[
 \mathcal S_{q,c}(0)=C_{q_2}(0)\,q_1^2\chi(q_1)\,\xi_{q_1}(0)\,S_c(0).
\]
By \eqref{eq:weightproperties} we have $\xi_{q_1}(0)=0$ if $q_1>1$,
while $\xi_1(0)=1$ by definition. Hence the terms with $u=0$
contribute
\begin{equation}\label{eq:zeroterm}
 c_Q\frac{B^2C_q(0)}{q^3}
 \sum_{\substack{c<r_0Q\\(c,q)=1}}\frac{S_c(0)}{c^4}\,I_c(0).
\end{equation}
The weight $w$ is bounded uniformly in $W$ and $y$, it is supported
in $\rho_1\le\lVert x\rVert\le\rho_2$, where $\abs{\nabla F}$ is
bounded below, and the part of the cone $F=0$ in this region has
finite surface measure. Hence $\sigma_\infty(w,F)\ll1$, and then
$I_r^*(0;w)\ll1$ for every $r>0$, by \eqref{eq:Izero} for $r<r_0$.

Set $X_0=Q^{1-\eta}$, so that $1\le X_0\le r_0Q$. For $c\le X_0$ we
have $c/Q\le Q^{-\eta}$, and \eqref{eq:Izero} with an integer
$j\ge2/\eta$ gives $I_c(0)=\sigma_\infty(w,F)+O(Q^{-2})$. By
Lemma~\ref{lem:badprime},
\[
 \sum_{\substack{c\le X_0\\(c,q)=1}}\frac{S_c(0)}{c^4}\,I_c(0)
 =\sigma_\infty(w,F)\sum_{\substack{c\le X_0\\(c,q)=1}}\frac{S_c(0)}{c^4}
 +O(Q^{-2}\log Q).
\]
By \eqref{eq:seriesfactor} and \eqref{eq:seriestail} with
$\eps=\eta$, the last sum over $c$ is
\[
 \S_F\prod_{p\mid q}\sigma_p^{-1}+O(q^\eta X_0^{-1+\eta}).
\]
For the remaining $c$ let $T(t)$ be the sum of $S_c(0)$ over the
integers $c\le t$ with $(c,q)=1$, and let
$\phi_0(t)=t^{-4}I_{t/Q}^*(0;w)$. Then $T(t)\ll t^3(tq)^\eta$ by
Lemma~\ref{lem:conemod}, and $\phi_0(t)\ll t^{-4}$. By the second
part of \eqref{eq:Izero},
\[
 \phi_0'(t)=-4t^{-5}I_r^*(0;w)
 +t^{-5}\,r\,\frac{\partial I_r^*(0;w)}{\partial r}\ll t^{-5},
 \qquad r=\frac tQ.
\]
Since $\phi_0(r_0Q)=0$, partial summation gives
\[
 \sum_{\substack{X_0<c<r_0Q\\(c,q)=1}}S_c(0)\,\phi_0(c)
 =-T(X_0)\phi_0(X_0)-\int_{X_0}^{r_0Q}T(t)\,\phi_0'(t)\,dt
 \ll q^\eta X_0^{-1+\eta}.
\]
Here $X_0^{-1+\eta}\le Q^{-1+2\eta}$. Collecting these estimates, we
find that the sum over $c$ in \eqref{eq:zeroterm} is
\[
 \sigma_\infty(w,F)\,\S_F\prod_{p\mid q}\sigma_p^{-1}
 +O(q^\eta Q^{-1+2\eta}).
\]
We multiply by $C_q(0)q^{-3}$, where
$\abs{C_q(0)}\le(2K)^{\omega(q)}q^2\ll q^{2+\eta}$ by
\eqref{eq:Cstrata}, and recall \eqref{eq:Sxi}. Thus
\eqref{eq:zeroterm} is
\[
 c_QB^2\big(\sigma_\infty(w,F)\,\S_\xi(q)
 +O(q^{-1+2\eta}Q^{-1+2\eta})\big).
\]
Finally $c_Q=1+O(Q^{-1})$ by Lemma~\ref{lem:delta}, and
$\S_\xi(q)\ll q^{-1+\eta}$ by \eqref{eq:Sxibound}. Since
$Q=B/\sqrt q$, the two error terms are
\[
 B^2Q^{-1}q^{-1+\eta}=Bq^{-1/2+\eta},\qquad
 B^2q^{-1+2\eta}Q^{-1+2\eta}=B^{1+2\eta}q^{-1/2+\eta},
\]
and both are at most $B(Bq)^{2\eta}$.
\end{proof}

\begin{lemma}\label{lem:annulus}
Let $\mathcal W$ be a set of weights as in
Lemma~\ref{lem:integrals}. For every $W\in\mathcal W$, every
squarefree $q\ge1$ with $(q,N)=1$ and every $B\ge2$,
\begin{equation}\label{eq:annulus}
 T_q(B;W)=B^2\sigma_\infty(W,F)\,\S_\xi(q)
 +O_{F,\mathcal W,E,K,\eps}\big((Bq)^{1+\eps}\big).
\end{equation}
\end{lemma}

\begin{proof}
We may assume that $\eps\le2$. Suppose first that $B\ge9$ and
$q\le B$, and let $\delta$ be as in Lemma~\ref{lem:integrals}. By
\eqref{eq:pieces},
\[
 T_q(B;W)=\delta^{-4}\int_{\R^4}T_q(B;W_y)\,dy,\qquad
 \sigma_\infty(W,F)=\delta^{-4}\int_{\R^4}\sigma_\infty(W_y,F)\,dy,
\]
where the integrands vanish unless $\lVert y\rVert\le\rho_2+2$; the
first identity is \eqref{eq:pieces} at each of the finitely many $x$
in \eqref{eq:Tw}, and in the second the order of the integrations
over $y$ and over the cone may be changed because $W$ is bounded and
$\abs{\nabla F}$ is bounded below on its support. If $F$ has no zero
in the support of $W_y$, then $T_q(B;W_y)=0$ and
$\sigma_\infty(W_y,F)=0$. For the other $y$,
Lemmas~\ref{lem:offcone}, \ref{lem:oncone} and~\ref{lem:zerofreq}
with $\eta=\eps/8$ give
\[
 T_q(B;W_y)=B^2\sigma_\infty(W_y,F)\,\S_\xi(q)+O\big((Bq)^{1+\eps}\big),
\]
with an implied constant that does not depend on $y$ or on $W$.
Integrating over $y$, we obtain \eqref{eq:annulus}.

In the remaining cases we have $q>B$ or $B<9$. The vectors $x$ in
\eqref{eq:Tw} satisfy $\lVert x\rVert\le\rho_2B$, and by
Lemma~\ref{lem:quadriccount} there are $O(B^{2+\eps})$ of them, so
that $T_q(B;W)\ll B^{2+\eps}$. Moreover
$\sigma_\infty(W,F)\ll1$, as in the proof of
Lemma~\ref{lem:zerofreq}, and
$B^2\sigma_\infty(W,F)\,\S_\xi(q)\ll B^2q^{-1+\eps}$ by
\eqref{eq:Sxibound}. Both bounds are $O((Bq)^{1+\eps})$ when $q>B$,
and also when $B<9$.
\end{proof}

\begin{proof}[Proof of Proposition~\ref{prop:analytic}]
It remains to remove the condition on the support of the weight. Let
$\kappa:\R\to[0,1]$ be an infinitely differentiable, nonincreasing
function with $\kappa(t)=1$ for $t\le1$ and $\kappa(t)=0$ for
$t\ge2$, and set
\[
 \omega_0(x)=\kappa(\lVert x\rVert)-\kappa(2\lVert x\rVert).
\]
Then $\omega_0$ is a weight supported in
$1/2\le\lVert x\rVert\le2$, it is nonnegative, and
\begin{equation}\label{eq:dyadic}
 \sum_{j=0}^J\omega_0(2^jx)=\kappa(\lVert x\rVert)-\kappa(2^{J+1}\lVert x\rVert)
 \qquad(J\ge0).
\end{equation}
Let $\rho\ge1$ be such that $W(x)=0$ for $\lVert x\rVert\ge\rho$, and
for $j\ge0$ set
\[
 W_j(x)=W(2^{-j}x)\,\omega_0(x/\rho),\qquad B_j=2^{-j}B.
\]
The weights $W_j$ are supported in $\rho/2\le\lVert x\rVert\le2\rho$,
and their partial derivatives of each order are bounded uniformly in
$j$. Thus Lemma~\ref{lem:annulus} applies to the set $\mathcal W$ of
the $W_j$, with implied constants depending only on $F$, $E$, $K$,
$W$ and $\eps$.

Since $W(x)\kappa(\lVert x\rVert/\rho)=W(x)$, the identity
\eqref{eq:dyadic} gives
\[
 \sum_{j=0}^JW_j(2^jx)=W(x)\big(1-\kappa(2^{J+1}\lVert x\rVert/\rho)\big),
\]
and the last factor is $1$ when $\lVert x\rVert\ge2^{-J}\rho$. Let
$J$ be the least integer with $2^J\ge\rho B$. Then
$W(x/B)=\sum_{j\le J}W_j(x/B_j)$ for every $x\in\Z^4$ other than the
origin, and therefore
\begin{equation}\label{eq:dyadicsum}
 T_q(B;W)=W(0)\,\xi_q(0)+\sum_{j=0}^JT_q(B_j;W_j).
\end{equation}
To the terms with $B_j\ge2$ we apply \eqref{eq:annulus}. The sum of
their error terms is $O((Bq)^{1+\eps})$. In the sum of the other
terms, a vector $x$ has the coefficient
$\xi_q(x)W(x/B)\sum_j\omega_0(2^jx/(\rho B))$, where $j$ runs over
the indices with $B_j<2$; this is at most $\max\abs W$ in absolute
value, by \eqref{eq:dyadic}, and it vanishes unless
$\lVert x\rVert<4\rho$. Hence the terms of \eqref{eq:dyadicsum} with
$B_j<2$ contribute $O_W(1)$, and so does the term $W(0)\xi_q(0)$.

It remains to consider the main terms. The substitution
$y=2^{-j}z$ maps the cone $F=0$ to itself, multiplies $dS$ by
$2^{-3j}$ and $\abs{\nabla F}$ by $2^{-j}$. Hence
\[
 B_j^2\sigma_\infty(W_j,F)=B^2\sigma_\infty(V_j,F),\qquad
 V_j(x)=W_j(2^jx)=W(x)\,\omega_0(2^jx/\rho).
\]
By \eqref{eq:dyadic} the sum of the $V_j(x)$ over all $j\ge0$ is
$W(x)$ for $x\ne0$, and its partial sums are at most $\abs{W(x)}$ in
absolute value. Since the integral defining $\sigma_\infty(W,F)$
converges absolutely, it follows that
\[
 \sum_{j\ge0}B_j^2\sigma_\infty(W_j,F)=B^2\sigma_\infty(W,F).
\]
The terms of this series with $B_j<2$ have the sum $O_{F,W}(1)$,
because $\sigma_\infty(W_j,F)$ is bounded uniformly in $j$. With
\eqref{eq:Sxibound} this gives \eqref{eq:analytic}.
\end{proof}

\subsection{The character sums}\label{sec:character-sums}

We deduce Theorems~\ref{thm:polynomial}, \ref{thm:diagonal}
and~\ref{thm:cyclic} from Proposition~\ref{prop:analytic} and the
results of Section~\ref{sec:fourier}. The following lemma contains
the first and the third of them.

\begin{lemma}\label{lem:charsum}
Let $n\ge2$, and let $G\in\Z[x_1,x_2,x_3,x_4]$ be a form of degree
$n$ satisfying \eqref{eq:smooth}. Let $\mathcal P$ be a set of primes
not dividing $N_G$, and for every $p\in\mathcal P$ let $\psi_p$ be a
nonprincipal character modulo $p$ such that $\psi_p^n$ is principal.
For $p\in\mathcal P$ let $a_{p,\psi}(G)$ be the sum
\eqref{eq:apsidef}, which does not depend on the coordinates chosen
for the points $P$ because $\psi_p^n$ is principal, and set
\begin{equation}\label{eq:localfactor}
 \beta_p=\frac{a_{p,\psi}(G)(p^2-1)}{p^2\,\#\Qc(\Fp)}.
\end{equation}
Then
\begin{equation}\label{eq:betabound}
 \abs{a_{p,\psi}(G)}\le(2n^2-4n+4)p,\qquad
 \beta_p=\frac{a_{p,\psi}(G)(p-1)(p-\chi(p))}{p^2(p^2-\chi(p))},\qquad
 \abs{\beta_p}\le\frac{2n^2-4n+4}p.
\end{equation}
Moreover, for every weight $W$, every squarefree $q\ge1$ all of whose
prime divisors lie in $\mathcal P$, and every $B\ge2$,
\begin{equation}\label{eq:charsum}
 \sum_{\substack{x\in\Z^4\\F(x)=0}}W(x/B)\prod_{p\mid q}\psi_p(G(x))
 =B^2\sigma_\infty(W,F)\,\S_F\prod_{p\mid q}\beta_p
 +O_{F,n,W,\eps}\big((Bq)^{1+\eps}\big).
\end{equation}
\end{lemma}

\begin{proof}
Let $p\in\mathcal P$. Since $G$ has good reduction at $p$, the first
bound in \eqref{eq:betabound} is Lemma~\ref{lem:katz}. By
Lemma~\ref{lem:gauss}(iii) and because $\chi(p)^2=1$,
\[
 (p-\chi(p))\,\#\Qc(\Fp)=(p-\chi(p))\big(p^2+(1+\chi(p))p+1\big)
 =(p+1)(p^2-\chi(p)),
\]
which gives the second expression for $\beta_p$. The third assertion
follows from the first two, because
$(p-1)(p-\chi(p))\le p^2-\chi(p)$.

By Lemma~\ref{lem:exceptional}(i) the integer $N_G$ is well defined.
We apply the results of Sections~\ref{sec:local}
and~\ref{sec:analytic} with $N=N_G$, with $E=E_{F,n}$ and with
$K=2n^2$. For $p\in\mathcal P$ we take $\xi_p(x)=\psi_p(G(x))$ and
the polynomial $D_p$ of Lemma~\ref{lem:exceptional}(ii); then
\eqref{eq:weightproperties} and \eqref{eq:Cstrata} hold by
Lemma~\ref{lem:fourier}. For the other primes $p\nmid N_G$ we take
$\xi_p=0$ and $D_p=1$; then $C_p(u)=0$ for every $u$, and
\eqref{eq:weightproperties} and \eqref{eq:Cstrata} hold as well. For
$q$ as in the lemma, $\xi_q(x)$ is the product of the
$\psi_p(G(x))$ over $p\mid q$, and Proposition~\ref{prop:analytic}
gives \eqref{eq:charsum} with $\S_\xi(q)$ in place of
$\S_F\prod_{p\mid q}\beta_p$. Now let $p\mid q$. We have
$C_p(0)=(p-1)a_{p,\psi}(G)$ by \eqref{eq:fourier}, and
$\sigma_p=(p^2-\chi(p))/(p(p-\chi(p)))$ by \eqref{eq:sigmagood}.
Hence
\[
 \frac{C_p(0)}{p^3\sigma_p}
 =\frac{a_{p,\psi}(G)(p-1)(p-\chi(p))}{p^2(p^2-\chi(p))}=\beta_p,
\]
and $\S_\xi(q)=\S_F\prod_{p\mid q}\beta_p$ by \eqref{eq:Sxi}.
\end{proof}

\begin{proof}[Proof of Theorem~\ref{thm:polynomial}]
The form $G$ has degree $n=2d$. By Lemma~\ref{lem:exceptional}(i)
the product in \eqref{eq:NG} is finite, and \eqref{eq:badheight}
holds. Every prime $p\nmid N_G$ is odd, so that $\chi_p$ is a
nonprincipal character modulo $p$ and $\chi_p^{2d}$ is principal. We
apply Lemma~\ref{lem:charsum} with the set $\mathcal P$ of all primes
not dividing $N_G$ and with $\psi_p=\chi_p$. Then
$a_{p,\psi}(G)=a_p(G)$, and the product of the $\chi_p(G(x))$ over
$p\mid q$ is $\chi_q(G(x))$. Thus \eqref{eq:charsum} is
\eqref{eq:mainthm}, with the two expressions of \eqref{eq:Sgeneral}
for $\S_G(q)$, and the first part of \eqref{eq:betabound} is the
bound for $a_p(G)$, because $2n^2-4n+4=8d^2-8d+4$. By the last part
of \eqref{eq:betabound},
\[
 \abs{\S_G(q)}\le\abs{\S_F}\frac{(8d^2-8d+4)^{\omega(q)}}q
 \ll_{F,d,\eps}q^{-1+\eps},
\]
and with \eqref{eq:mainthm} this gives the bound for $T_{q,G}(B)$ in
\eqref{eq:mainbound}.
\end{proof}

\begin{proof}[Proof of Theorem~\ref{thm:cyclic}]
The form $G$ has degree $n=rd$, and the first assertion is
Lemma~\ref{lem:exceptional}(i). We apply Lemma~\ref{lem:charsum} with
the set $\mathcal P$ of the primes $p\equiv1\imod r$ not dividing
$N_G$ and with the given characters $\psi_p$, which are nonprincipal
because $r\ge3$, and whose $n$th powers are principal. Then
\eqref{eq:charsum} is \eqref{eq:cyclicmain} with
\eqref{eq:cyclicfactor}, and the first part of \eqref{eq:betabound}
is the bound for $a_{p,\psi}(G)$, because
$2n^2-4n+4=2(1+(rd-1)^2)$.
\end{proof}

\begin{proof}[Proof of Theorem~\ref{thm:diagonal}]
The first two assertions are Lemma~\ref{lem:diagbad} and
\eqref{eq:diagtrace}, which was proved in Lemma~\ref{lem:pencil}. By
Lemma~\ref{lem:diagbad} the form $G$ satisfies \eqref{eq:smooth},
and the condition $(q,N_G)=1$ is $(q,2\alpha M)=1$. Hence
Theorem~\ref{thm:polynomial} applies with $d=1$, and inserting
\eqref{eq:diagtrace} into the second expression of
\eqref{eq:Sgeneral} gives \eqref{eq:diagmain}. We also verify the
identity $\kappa_1\kappa_2\kappa_3\kappa_4=\alpha M^2$ used after the
theorem: the product of the four factors $-a_i$ is $\alpha$, and for
each pair $i<j$ the product of the $\kappa_i$ contains the two
factors $a_ib_j-a_jb_i$ and $a_jb_i-a_ib_j$, whose product over the
six pairs is $M^2$.
\end{proof}

For $p\nmid N_G$ the number $\beta_p$ of \eqref{eq:localfactor},
formed with $\psi_p=\chi_p$, is the factor at $p$ in
\eqref{eq:Sgeneral}. When it is subtracted from each character, the
main term disappears.

\begin{corollary}\label{cor:centered}
Let $d\ge1$, and let $G$ be a form of degree $2d$ satisfying
\eqref{eq:smooth}. For every weight $W$, every squarefree $q>1$ with
$(q,N_G)=1$, and every $B\ge2$,
\begin{equation}\label{eq:centered}
 \sum_{\substack{x\in\Z^4\\F(x)=0}}W(x/B)
 \prod_{p\mid q}\big(\chi_p(G(x))-\beta_p\big)
 \ll_{F,d,W,\eps}(Bq)^{1+\eps}.
\end{equation}
\end{corollary}

\begin{proof}
For a divisor $s$ of $q$ write $\beta_s$ for the product of the
$\beta_p$ over $p\mid s$. Expanding the product, we see that the left
side of \eqref{eq:centered} is
\[
 \sum_{st=q}\mu(s)\beta_s\,T_{t,G}(B)
 =B^2\sigma_\infty(W,F)\,\S_F\sum_{st=q}\mu(s)\beta_s\beta_t
 +O\bigg(\sum_{st=q}\abs{\beta_s}(Bt)^{1+\eps}\bigg),
\]
by Theorem~\ref{thm:polynomial}. Here $\beta_s\beta_t=\beta_q$, and
the sum of $\mu(s)$ over the divisors $s$ of $q$ vanishes because
$q>1$. In the error term $(Bt)^{1+\eps}\le(Bq)^{1+\eps}s^{-1}$, and
by \eqref{eq:betabound}
\[
 \sum_{s\mid q}\frac{\abs{\beta_s}}s
 \le\prod_p\Big(1+\frac{8d^2-8d+4}{p^2}\Big)\ll_d1.\qedhere
\]
\end{proof}

\subsection{Square and power values}\label{sec:sieve}

Theorem~\ref{thm:uniformcount} and Corollary~\ref{cor:powercount}
are the cases $r=2$, $n=2d$ and $r\ge3$, $n=rd$ of the following
lemma. Its proof is the square sieve of Heath-Brown \cite{HB2}, and
for $r\ge3$ the analogue for $r$th powers used in
\cite[Section~2]{HBP}, with
Lemma~\ref{lem:charsum} at the moduli $q=pp'$.

\begin{lemma}\label{lem:sieve}
Let $n\ge2$, let $r\ge2$ be a divisor of $n$, and let
$G\in\Z[x_1,x_2,x_3,x_4]$ be a form of degree $n$ satisfying
\eqref{eq:smooth}. Let $M_{G,r}(B)$ be the number of $x\in\Z^4$ with
$\lVert x\rVert_\infty\le B$ and $F(x)=0$ for which $G(x)=k^r$ for
some integer $k\ne0$. Then, for $B\ge2$,
\begin{equation}\label{eq:sieve}
 M_{G,r}(B)\ll_{F,n,\eps}B^{5/3+\eps}+B^{1+\eps}\log^2(2H_GB).
\end{equation}
\end{lemma}

\begin{proof}
Write $\mathcal L=\log(2H_GB)$, so that $\mathcal L>1$. The vectors
$x$ counted by $M_{G,r}(B)$ satisfy $\lVert x\rVert\le2B$ and
$F(x)=0$, and Lemma~\ref{lem:quadriccount} gives
$M_{G,r}(B)\ll_{F,\eps}B^{2+\eps}$. This proves \eqref{eq:sieve} if
$\mathcal L\ge B^{1/2}$, and we assume from now on that
$\mathcal L<B^{1/2}$.

In this proof $c_1,c_2,c_3$ and $z_1\ge2$ are positive numbers
depending only on $F$ and $n$. By the prime number theorem for arithmetic
progressions \cite[Corollary~11.21]{MV}, for $z\ge z_1$ the interval
$(z,2z]$ contains at least $c_1z/\log z$ primes $p\equiv1\imod r$. By
Lemma~\ref{lem:exceptional}(i) we have $\log N_G\le c_2\mathcal L$.
For $x\in\Z^4$ with $\lVert x\rVert_\infty\le B$ and $G(x)\ne0$ we
have $\abs{G(x)}\le\binom{n+3}3H_GB^n$, and therefore
$\log\abs{G(x)}\le c_3\mathcal L$. Set
\[
 z=c_4(B^{1/3}+\mathcal L),\qquad
 c_4=\max\{z_1,\,2c_2/c_1,\,4c_3/c_1\},
\]
and let $\mathcal P$ be the set of the primes $p\equiv1\imod r$ with
$z<p\le2z$ and $p\nmid N_G$. At most $\log N_G/\log z$ primes $p>z$
divide $N_G$. Since $z\ge z_1$ and $c_2\mathcal L\le c_1z/2$, it
follows that
\begin{equation}\label{eq:sieveprimes}
 \#\mathcal P\ge\frac{c_1z}{2\log z}.
\end{equation}
For every $p\in\mathcal P$ let $\psi_p$ be a character modulo $p$ of
exact order $r$; for $r=2$ this is $\chi_p$.

Let $x$ be one of the vectors counted by $M_{G,r}(B)$, so that
$G(x)=k^r$ with an integer $k\ne0$. For $p\in\mathcal P$ we have
$\psi_p(G(x))=\psi_p(k)^r=1$ if $p\nmid G(x)$, and $\psi_p(G(x))=0$
otherwise. The number of $p\in\mathcal P$ dividing $G(x)$ is at most
$\log\abs{G(x)}/\log z\le c_3\mathcal L/\log z$, which is at most
$\#\mathcal P/2$ by \eqref{eq:sieveprimes} and the choice of $c_4$.
Hence
\[
 \sum_{p\in\mathcal P}\psi_p(G(x))\ge\frac{\#\mathcal P}2.
\]
Let $\kappa$ be as in the proof of
Proposition~\ref{prop:analytic}, and let
$W(x)=\prod_{i=1}^4\kappa(\abs{x_i})$. This is a nonnegative weight
with $W(x)=1$ for $\lVert x\rVert_\infty\le1$ and $W(x)=0$ for
$\lVert x\rVert_\infty\ge2$. It follows that
\begin{equation}\label{eq:sieveineq}
\begin{aligned}
 M_{G,r}(B)&\le\frac4{(\#\mathcal P)^2}
 \sum_{\substack{x\in\Z^4\\F(x)=0}}W(x/B)\,
 \bigg|\sum_{p\in\mathcal P}\psi_p(G(x))\bigg|^2\\
 &=\frac4{(\#\mathcal P)^2}\sum_{p,p'\in\mathcal P}\
 \sum_{\substack{x\in\Z^4\\F(x)=0}}W(x/B)\,
 \psi_p(G(x))\overline{\psi_{p'}(G(x))}.
\end{aligned}
\end{equation}

Each of the $\#\mathcal P$ terms with $p=p'$ is at most the sum of
$W(x/B)$ over the zeros $x$ of $F$, which is
$O_{F,\eps}(B^{2+\eps})$ by Lemma~\ref{lem:quadriccount}. Let $p\ne p'$. We apply
Lemma~\ref{lem:charsum} with the set $\{p,p'\}$, with the characters
$\psi_p$ and $\overline{\psi_{p'}}$, and with $q=pp'$. By
\eqref{eq:charsum} and the last part of \eqref{eq:betabound},
\[
 \sum_{\substack{x\in\Z^4\\F(x)=0}}W(x/B)\,
 \psi_p(G(x))\overline{\psi_{p'}(G(x))}
 \ll_{F,n,\eps}\frac{B^2}{pp'}+(Bpp')^{1+\eps}
 \le\frac{B^2}{z^2}+(4Bz^2)^{1+\eps}.
\]
Inserting these bounds and \eqref{eq:sieveprimes} into
\eqref{eq:sieveineq}, we obtain
\[
 M_{G,r}(B)\ll_{F,n,\eps}\frac{B^{2+\eps}\log z}z
 +\frac{B^2}{z^2}+(Bz^2)^{1+\eps}.
\]
Since $\mathcal L<B^{1/2}$, we have $B^{1/3}\le z\ll_{F,n}B^{1/2}$,
so that $\log z\ll_{F,n}\log B$ and
$(Bz^2)^\eps\ll_{F,n}B^{2\eps}$. Moreover
$z^2\ll_{F,n}B^{2/3}+\mathcal L^2$. Hence
\[
 M_{G,r}(B)\ll_{F,n,\eps}B^{5/3+2\eps}+B^{4/3}
 +B^{2\eps}(B^{5/3}+B\mathcal L^2),
\]
which gives \eqref{eq:sieve} with $2\eps$ in place of $\eps$.
\end{proof}

\subsection{The ranges of the modulus}\label{sec:relative}

We prove the statements made at the end of
Section~\ref{sec:results}. They rest on lower bounds for the local
factors $\beta_p$ of \eqref{eq:localfactor} when these do not vanish.
For a quadratic form $G$ the sum $a_p(G)$ is divisible by $p$, as in
\eqref{eq:diagtrace}.

\begin{lemma}\label{lem:divisible}
Let $G$ be a quadratic form satisfying \eqref{eq:smooth}, and let
$p\nmid N_G$ be a prime. Then $a_p(G)$ is an integer divisible by
$p$, and $\abs{a_p(G)}\le4p$.
\end{lemma}

\begin{proof}
The bound is Lemma~\ref{lem:katz} with $n=2$ and $\psi=\chi_p$. Let
$C_p(0)$ be the sum \eqref{eq:Cpsi} with $\psi=\chi_p$ and $u=0$. For
$m\in\Fp$ we have $\chi_p(m)\equiv m^{(p-1)/2}\imod p$, and
$1-m^{p-1}$ is $1$ if $m=0$ and is $0$ otherwise. Hence
\[
 C_p(0)\equiv\sum_{x\in\Fp^4}\big(1-F(x)^{p-1}\big)G(x)^{(p-1)/2}\imod p.
\]
The polynomial on the right has degree at most $3(p-1)$, so that each
of its monomials $x_1^{k_1}x_2^{k_2}x_3^{k_3}x_4^{k_4}$ has an
exponent $k_i<p-1$. For $0\le k<p-1$ the sum of $y^k$ over
$y\in\Fp$ vanishes in $\Fp$: this is clear for $k=0$, and for $k>0$
the sum is unchanged when it is multiplied by $g^k\ne1$, where $g$ is
a generator of the multiplicative group of $\Fp$. Hence
$C_p(0)\equiv0\imod p$. Since $C_p(0)=(p-1)a_p(G)$ by
\eqref{eq:fourier}, the lemma follows.
\end{proof}

\begin{lemma}\label{lem:lowerbeta}
Let $n$, $G$, $\mathcal P$ and the characters $\psi_p$ be as in
Lemma~\ref{lem:charsum}, and let $p\in\mathcal P$ be a prime with
$a_{p,\psi}(G)\ne0$.
\begin{enumerate}
\item[\textup{(i)}] If $\psi_p=\chi_p$, then
$\abs{\beta_p}\ge1/(2p^2)$, and $\abs{\beta_p}\ge1/(2p)$ if $n=2$.
\item[\textup{(ii)}] If $\psi_p$ has exact order $r\ge3$, then
\[
 \abs{\beta_p}\ge\frac1{2p^2}\big((2n^2-4n+4)p\big)^{1-\varphi(r)/2}.
\]
\end{enumerate}
\end{lemma}

\begin{proof}
We have $(p-1)(p-\chi(p))\ge(p^2-\chi(p))/2$, because $p\ge3$. By
the second part of \eqref{eq:betabound}, therefore,
$\abs{\beta_p}\ge\abs{a_{p,\psi}(G)}/(2p^2)$.

(i) If $\psi_p=\chi_p$, then $a_{p,\psi}(G)=a_p(G)$ is an integer, so
that $\abs{a_p(G)}\ge1$. If $n=2$, then $\abs{a_p(G)}\ge p$ by
Lemma~\ref{lem:divisible}.

(ii) Let $\zeta$ be a primitive $r$th root of unity. The values of
$\psi_p$ are $0$ and powers of $\zeta$, so that
$a=a_{p,\psi}(G)$ belongs to $\Z[\zeta]$. The automorphisms of the
field $\Q(\zeta)$ are the $\varphi(r)$ maps $\sigma_t$ with
$\sigma_t(\zeta)=\zeta^t$, where $t$ runs over the reduced residues
modulo $r$ \cite[Chapter~VI, Theorem~3.1]{La}. Since
$\sigma_t(\psi_p(m))=\psi_p^t(m)$ for every integer $m$, the number
$\sigma_t(a)$ is the sum \eqref{eq:apsidef} formed with the character
$\psi_p^t$. This character has exact order $r$, and
Lemma~\ref{lem:katz} gives $\abs{\sigma_t(a)}\le(2n^2-4n+4)p$. The
product of the $\sigma_t(a)$ is the norm of $a$, a rational integer,
and it is not zero because $a\ne0$. Moreover
$\abs{\sigma_{-1}(a)}=\abs a$, because $\sigma_{-1}$ is the complex
conjugation, and $-1\not\equiv1\imod r$. Hence
\[
 1\le\prod_t\abs{\sigma_t(a)}\le\abs a^2\big((2n^2-4n+4)p\big)^{\varphi(r)-2},
\]
which gives the assertion.
\end{proof}

\begin{corollary}\label{cor:relative}
Let $n$, $G$, $\mathcal P$ and the characters $\psi_p$ be as in
Lemma~\ref{lem:charsum}. Let $W$ be a weight with
$\sigma_\infty(W,F)\,\S_F\ne0$, and let $q\ge1$ be a squarefree
integer, all of whose prime divisors lie in $\mathcal P$, such that
$a_{p,\psi}(G)\ne0$ for every $p\mid q$. Suppose that either
$\psi_p=\chi_p$ for every $p\mid q$, and set $\gamma=2$ if $n=2$ and
$\gamma=3$ if $n>2$; or that, for some $r\ge3$, the character
$\psi_p$ has exact order $r$ for every $p\mid q$, and set
$\gamma=2+\varphi(r)/2$. Then, for $B\ge2$,
\begin{equation}\label{eq:relative}
 \sum_{\substack{x\in\Z^4\\F(x)=0}}W(x/B)\prod_{p\mid q}\psi_p(G(x))
 =B^2\sigma_\infty(W,F)\,\S_F\prod_{p\mid q}\beta_p\,
 \bigg(1+O_{F,n,W,\eps}\Big(\frac{q^\gamma}B(Bq)^\eps\Big)\bigg).
\end{equation}
\end{corollary}

\begin{proof}
By Lemma~\ref{lem:lowerbeta} there is a number $c\ge1$, depending
only on $n$, such that $\abs{\beta_p}\ge(cp^{\gamma-1})^{-1}$ for
every $p\mid q$; in case (ii) we use that $r\le n$. Hence the product
of the $\beta_p$ over $p\mid q$ is at least
$c^{-\omega(q)}q^{1-\gamma}\gg_{n,\eps}q^{1-\gamma-\eps}$ in absolute
value, and \eqref{eq:relative} follows on dividing the error term of
\eqref{eq:charsum} by the main term.
\end{proof}

The error term in \eqref{eq:relative} tends to zero uniformly for
$q\le B^{1/\gamma-\eta}$, for every fixed $\eta>0$. With
$\psi_p=\chi_p$ this is the range $q\le B^{1/2-\eta}$ for a quadratic
form $G$ and $q\le B^{1/3-\eta}$ in general, and for characters of
order $r\ge3$ it is the range $q\le B^{2/(\varphi(r)+4)-\eta}$, as
stated in Section~\ref{sec:results}.

\section{Comparison with the determinant method}\label{sec:fixed}

For a form $G$ of degree $r$, the number $M_{G,r}(B)$ of
Lemma~\ref{lem:sieve} can also be bounded by the determinant method,
through a theorem of Salberger \cite{S} on complete intersections in
$\PP^4$. For a quasi-projective variety $U\subseteq\PP^4$ defined
over $\Q$ and $B\ge1$, let $N(U;B)$ be the number of points of
$U(\Q)$ that are represented by a vector
$(x_1,\ldots,x_5)\in\Z^5$ with coprime coordinates and
$\max_i\abs{x_i}\le B$, as in \cite{S}.

\begin{proposition}\label{prop:salberger}
Let $r\ge2$, and let $G\in\Z[x_1,x_2,x_3,x_4]$ be a form of degree
$r$ satisfying \eqref{eq:smooth}. Then, for $B\ge1$,
\begin{equation}\label{eq:salberger}
 M_{G,r}(B)\ll_{r,\eps}(H_G^{1/r}B)^{m_r+\eps},
 \qquad m_r=\max\Big\{1,\frac3{\sqrt{2r}}\Big\}.
\end{equation}
The implied constant depends only on $r$ and $\eps$.
\end{proposition}

\begin{proof}
Let $X\subset\PP^4$ be the variety defined by the equations
\[
 F(x)=0,\qquad y^r=G(x)
\]
in the coordinates $(x_1:x_2:x_3:x_4:y)$. We show that $X$ is a
nonsingular complete intersection of hypersurfaces of degrees $2$
and $r$. The gradients of the two equations at a point $(x:y)$ of
$X$ over $\overline\Q$ are $(\nabla F(x),0)$ and
$(-\nabla G(x),ry^{r-1})$. If $y\ne0$, they are linearly dependent
only if $\nabla F(x)=2Ax$ vanishes, that is, only if $x=0$; but then
$y^r=G(0)=0$. If $y=0$, then $x\ne0$ and $F(x)=G(x)=0$, and the
gradients are linearly independent by \eqref{eq:smooth}.

By \cite[Theorem~0.6]{S}, with $d_1=2$ and $d_2=r$, the complement
$U$ of the union of the lines on $X$ satisfies
$N(U;B)\ll_{r,\eps}B^{3/\sqrt{2r}+\eps}+B$, where the implied
constant depends only on $r$ and $\eps$. The number of lines on $X$
is bounded in terms of $r$; see \cite[p.~1094]{S}. We claim that no
line on $X$ is defined over $\Q$. Such a line does not contain the
point $(0:0:0:0:1)$, which is not on $X$, and its projection to the
coordinates $x$ is therefore a line in $\PP^3$, defined over $\Q$, on
which $F$ vanishes. Thus $F$ vanishes on a subspace of $\Q^4$ of
dimension $2$. Let $T$ be an invertible rational matrix whose first
two columns form a basis of this subspace. Then
\[
 T^{\mathsf T}AT=\begin{pmatrix}0&M_1\\M_1^{\mathsf T}&M_2\end{pmatrix}
\]
with $2\times2$ blocks, and $\alpha\det(T)^2=\det(M_1)^2$, which is
impossible because $\alpha$ is not a square. A line on $X$ that is
not defined over $\Q$ contains at most one rational point, since two
rational points span a line defined over $\Q$. Hence
\begin{equation}\label{eq:salbergerX}
 N(X;B)\ll_{r,\eps}B^{m_r+\eps}\qquad(B\ge1).
\end{equation}

Let $x$ be one of the vectors counted by $M_{G,r}(B)$, so that
$G(x)=k^r$ with an integer $k\ne0$, and let $g$ be the greatest
common divisor of its coordinates. Then $x'=x/g$ has coprime
coordinates, $\lVert x'\rVert_\infty\le B/g$ and $F(x')=0$, and
$G(x')=(k/g)^r$ is an integer, so that $k'=k/g$ is an integer. The
point $(x':k')$ lies on $X$. A point of $X(\Q)$ arises in this way
from at most two vectors $x'$, which differ in sign. Since $G$
has $\binom{r+3}3$ coefficients,
$\abs{k'}^r\le\binom{r+3}3H_G(B/g)^r$. Hence, for each $g\le B$, the
number of vectors $x$ with this value of $g$ is at most
$2N(X;c_rH_G^{1/r}B/g)$, where $c_r=\binom{r+3}3^{1/r}$, and
\eqref{eq:salbergerX} gives
\[
 M_{G,r}(B)\ll_{r,\eps}\sum_{g\le B}\bigg(\frac{H_G^{1/r}B}g\bigg)^{m_r+\eps}
 \ll(H_G^{1/r}B)^{m_r+\eps}.\qedhere
\]
\end{proof}

For $r=2$ the bound \eqref{eq:salberger} is
$M_G(B)\ll_\eps(H_G^{1/2}B)^{3/2+\eps}$. We compare
\eqref{eq:salberger} with \eqref{eq:uniformcount} and
\eqref{eq:powercount} for a form $G$ of degree $r$. The powers of
$B$ agree when $H_G=B^{\theta_r}$, where
\[
 \theta_r=r\Big(\frac5{3m_r}-1\Big),\qquad\text{so that}\qquad
 \theta_2=\frac29,\quad\theta_3=1.08\ldots,\quad\theta_4=2.28\ldots,
 \quad\theta_r=\frac{2r}3\quad(r\ge5).
\]
For every fixed $\eta>0$, \eqref{eq:salberger} gives a smaller power
of $B$ when $H_G\le B^{\theta_r-\eta}$, and \eqref{eq:uniformcount}
and \eqref{eq:powercount} give a smaller power when
$H_G\ge B^{\theta_r+\eta}$. The dependence of
\eqref{eq:uniformcount} and \eqref{eq:powercount} on the
coefficients of $G$ is logarithmic, and these bounds hold for forms
$G$ of every degree divisible by $r$.

\section{Conclusion and future work}\label{sec:conclusion}

We have proved an asymptotic formula for the sums $T_{q,G}(B)$, and
for the corresponding sums with characters of order $r\ge3$, in which
the error term $(Bq)^{1+\eps}$ is uniform in the modulus $q$ and in
the coefficients of $G$, and the local factors are explicit. As an
application we have bounded the number of zeros of $F$ at which $G$
is a nonzero square or $r$th power.

\subsection{The error term}\label{sec:conclusion-error}

In Section~\ref{sec:analytic} the three kinds of frequencies are
treated in Lemmas~\ref{lem:offcone}, \ref{lem:oncone}
and~\ref{lem:zerofreq}, for weights supported in an annulus and for
$q\le B$. The frequency $u=0$ gives the
main term with an error $O(B(Bq)^\eps)$, and the frequencies $u\ne0$
with $F^*(u)=0$ contribute $O(B\sqrt q\,(Bq)^\eps)$. The term
$(Bq)^{1+\eps}$ comes from the frequencies with $F^*(u)\ne0$ alone,
for which the sums $\mathcal S_{q,c}(u)$ are bounded in absolute
value. We do not know whether the sum over these frequencies has
cancellation. An error term smaller than $Bq$ would extend the ranges
of $q$ at the end of Section~\ref{sec:results}.

\subsection{The exponent \texorpdfstring{$5/3$}{5/3}}\label{sec:conclusion-exponent}

The exponent $5/3$ in \eqref{eq:uniformcount} and
\eqref{eq:powercount} is that of \cite[Theorem~1.3]{M}. The proof
of Lemma~\ref{lem:sieve} uses primes $p$ of size
$B^{1/3}+\log(2H_GB)$. When $\log(2H_GB)\le B^{1/3}$, the terms with
$p=p'$ and the error term of \eqref{eq:charsum} at the moduli $q=pp'$
contribute the same power of $B$. Suppose that, for some
$1/2\le\theta<1$, the error term of \eqref{eq:charsum} were
$O(B^{1+\eps}q^{\theta+\eps})$ at the moduli $q=pp'$, with the same
uniformity. Then the same proof, with primes of size
$B^{1/(1+2\theta)}+\log(2H_GB)$, would give
\[
 M_{G,r}(B)\ll B^{2-1/(1+2\theta)+\eps}+B^{1+\eps}\log^2(2H_GB).
\]
For $\theta=1/2$, which corresponds to the contribution of
Lemma~\ref{lem:oncone}, the first exponent is $3/2$.

\subsection{The coefficients of \texorpdfstring{$G$}{G}}\label{sec:conclusion-coefficients}

The bounds \eqref{eq:uniformcount} and \eqref{eq:powercount} depend
on the coefficients of $G$ only through the term
$B^{1+\eps}\log^2(2H_GB)$. For a form $G$ of degree $r$ the bound
\eqref{eq:salberger} of the determinant method has a smaller power of
$B$ and a power of $H_G$; the two are compared at the end of
Section~\ref{sec:fixed}. We do not know whether $M_G(B)$ and
$M_{G,r}(B)$ have bounds with a power of $B$ smaller than $5/3$ in
which $H_G$ enters only through a power of $\log H_G$.

\subsection{Other cases}\label{sec:conclusion-cases}

Proposition~\ref{prop:analytic} is stated for functions $\xi_p$
satisfying \eqref{eq:weightproperties} and \eqref{eq:Cstrata}, and
the form $G$ enters only in Section~\ref{sec:fourier}. The
proposition therefore applies to other functions $\xi_p$ on the zeros
of $F$ modulo $p$, once \eqref{eq:Cstrata} is verified for them. The moduli
in this paper are squarefree and prime to $N_G$; we have not
considered moduli divisible by squares or by primes of bad reduction.
We have also assumed that $\alpha$ is not a square. This is used
when Lemma~\ref{lem:quadriccount} is applied to $F$ and to $F^*$, in
Lemmas~\ref{lem:conemod} and~\ref{lem:singseries}, and in
Section~\ref{sec:fixed}. When $\alpha$ is a square, the count
\eqref{eq:HB} is the sum of a constant multiple of $B^2\log B$, a
constant multiple of $B^2$ and $O(B^{3/2+\eps})$
\cite[Theorem~7]{HB}.

\section*{Acknowledgement}

The authors conceived this paper and determined the overall strategy
of the proofs. In developing and revising the arguments, they worked
with OpenAI's GPT-6 Astra, in particular on the treatment of the
exceptional frequencies in Section~\ref{sec:average}; the derivation
of the Fourier estimate \eqref{eq:fourier} and the construction of
its exceptional set in Lemma~\ref{lem:exceptional}, with bounds
uniform in the coefficients of $G$; the derivation of the uniform
estimate \eqref{eq:analytic} in Proposition~\ref{prop:analytic}; and
the comparison with the determinant method in
Proposition~\ref{prop:salberger}. It also assisted with the
translation of mathematical statements into Lean~\cite{NRZ}.
Anthropic's Claude (Fable~5.2) was used for adversarial audits and
to rewrite the authors' draft, which included
checking the arguments and citations against the sources and
reorganizing several proofs, notably those of
Lemma~\ref{lem:divisible} and Proposition~\ref{prop:salberger}.

The authors directed the work, determined the final content of the
paper, and take full responsibility for its mathematical statements,
proofs, references and exposition.

\begin{thebibliography}{BCLP}

\bibitem[BCLP]{BCLP}
A.~Bucur, A.~C. Cojocaru, M.~N. Lal\'in, and L.~B. Pierce,
\emph{Geometric generalizations of the square sieve, with an application to cyclic covers},
with an appendix by J.~Rabinoff,
Mathematika \textbf{69} (2023), no.~1, 106-154.

\bibitem[BKW]{BKW}
D.~Bonolis, E.~Kowalski, and K.~Woo,
\emph{Stratification theorems for exponential sums in families},
arXiv:2506.18299v3, 2026.

\bibitem[BL]{BL}
T.~D. Browning and D.~Loughran,
\emph{Sieving rational points on varieties},
Trans. Amer. Math. Soc. \textbf{371} (2019), no.~8, 5757-5785.

\bibitem[BM]{BM}
T.~D. Browning and R.~Munshi,
\emph{Rational points on singular intersections of quadrics},
Compos. Math. \textbf{149} (2013), no.~9, 1457-1494.

\bibitem[BP]{BP}
D.~Bonolis and L.~B. Pierce,
\emph{Application of a polynomial sieve: beyond separation of variables},
Algebra Number Theory \textbf{18} (2024), no.~8, 1515-1556;
correction in arXiv:2209.02494v3, 2026.

\bibitem[BPW]{BPW}
D.~Bonolis, L.~B. Pierce, and K.~Woo,
\emph{Counting integral points in thin sets of type II: singularities, sieves, and stratification},
arXiv:2505.11226v1, 2025.

\bibitem[DFI]{DFI}
W.~Duke, J.~B. Friedlander, and H.~Iwaniec,
\emph{Bounds for automorphic $L$-functions},
Invent. Math. \textbf{112} (1993), no.~1, 1-8.

\bibitem[G]{G}
J.~R. Getz,
\emph{Secondary terms in asymptotics for the number of zeros of quadratic forms over number fields},
J. Lond. Math. Soc. (2) \textbf{98} (2018), no.~2, 275-305.

\bibitem[GSS]{GSS}
R.~Guria, P.~Sharma, and S.~Singh,
\emph{On quadric surfaces weighted by arbitrary coefficients},
arXiv:2610.12944v1, 2026.

\bibitem[Ha]{Ha}
R.~Hartshorne,
\emph{Algebraic geometry},
Grad. Texts in Math., vol.~52, Springer, New York, 1977.

\bibitem[HB]{HB}
D.~R. Heath-Brown,
\emph{A new form of the circle method, and its application to quadratic forms},
J. Reine Angew. Math. \textbf{481} (1996), 149-206.

\bibitem[HB2]{HB2}
D.~R. Heath-Brown,
\emph{The square sieve and consecutive square-free numbers},
Math. Ann. \textbf{266} (1984), no.~3, 251-259.

\bibitem[HBP]{HBP}
D.~R. Heath-Brown and L.~B. Pierce,
\emph{Counting rational points on smooth cyclic covers},
J. Number Theory \textbf{132} (2012), no.~8, 1741-1757.

\bibitem[HSS]{HSS}
Z.~Huang, D.~Schindler, and A.~Shute,
\emph{Quantitative strong approximation for quaternary quadratic forms},
arXiv:2501.16766v1, 2025.

\bibitem[IK]{IK}
H.~Iwaniec and E.~Kowalski,
\emph{Analytic number theory},
Amer. Math. Soc. Colloq. Publ., vol.~53, Amer. Math. Soc., Providence, RI, 2004.

\bibitem[IM]{IM}
H.~Iwaniec and R.~Munshi,
\emph{The circle method and pairs of quadratic forms},
J. Th\'eor. Nombres Bordeaux \textbf{22} (2010), no.~2, 403-419.

\bibitem[K]{K}
N.~M. Katz,
\emph{Estimates for nonsingular multiplicative character sums},
Int. Math. Res. Not. \textbf{2002}, no.~7, 333-349.

\bibitem[L]{L}
S.~Lindqvist,
\emph{Weak approximation results for quadratic forms in four variables},
arXiv:1704.00502v1, 2017.

\bibitem[La]{La}
S.~Lang,
\emph{Algebra},
revised third edition, Grad. Texts in Math., vol.~211, Springer, New York, 2002.

\bibitem[M]{M}
K.~Mallesham,
\emph{On pairs of quadratic forms in five variables},
Int. J. Number Theory \textbf{17} (2021), no.~9, 2055-2070.

\bibitem[MV]{MV}
H.~L. Montgomery and R.~C. Vaughan,
\emph{Multiplicative number theory. I. Classical theory},
Cambridge Stud. Adv. Math., vol.~97, Cambridge Univ. Press, Cambridge, 2007.

\bibitem[NRZ]{NRZ}
H.~Nath, N.~Robles, and A.~Zaharescu,
\emph{Lean files for ``Character sums over zeros of quaternary quadratic forms''},
GitHub repository, \url{https://github.com/nicolasrobles/quaternary-quadratic-forms}, 2026.

\bibitem[S]{S}
P.~Salberger,
\emph{Counting rational points on projective varieties},
Proc. Lond. Math. Soc. (3) \textbf{126} (2023), no.~4, 1092-1133.

\end{thebibliography}

\end{document}
